%% (c) 2026 Brayden Sanders / 7Site LLC, Arkansas. All rights reserved.
%% For human use only. No commercial or government use without written agreement from 7Site LLC.

\documentclass[11pt]{article}
\usepackage{amsmath,amssymb,amsthm,mathtools}
\usepackage[margin=1in]{geometry}
\usepackage{hyperref}
\usepackage{enumitem}
\usepackage{booktabs}

\newtheorem{theorem}{Theorem}[section]
\newtheorem{lemma}[theorem]{Lemma}
\newtheorem{proposition}[theorem]{Proposition}
\newtheorem{corollary}[theorem]{Corollary}
\newtheorem{conjecture}[theorem]{Conjecture}
\theoremstyle{definition}
\newtheorem{definition}[theorem]{Definition}
\newtheorem{hypothesis}[theorem]{Hypothesis}
\newtheorem{example}[theorem]{Example}
\theoremstyle{remark}
\newtheorem{remark}[theorem]{Remark}

\DeclareMathOperator{\poly}{poly}
\DeclareMathOperator{\TV}{TV}
\newcommand{\calC}{\mathcal{C}}
\newcommand{\calD}{\mathcal{D}}
\newcommand{\calS}{\mathcal{S}}
\newcommand{\calF}{\mathcal{F}}
\newcommand{\Ppoly}{\mathrm{P/poly}}
\newcommand{\ACz}{\mathrm{AC}^0}
\newcommand{\NCone}{\mathrm{NC}^1}
\newcommand{\TCz}{\mathrm{TC}^0}
\newcommand{\SAT}{\mathrm{SAT}}
\newcommand{\threeSAT}{\mathrm{3\text{-}SAT}}
\newcommand{\dSAT}{\delta_{\mathrm{SAT}}}

\title{Irreducible Logical Entropy and the TIG Phantom Tile:\\
A Structural Analysis of the P--NP Gap\\[6pt]
{\large Coherence Field v1.0 --- February 2026}}
\author{Brayden Sanders / 7Site LLC\\[2pt]
{\normalsize Delta Signature Hash: \texttt{4b5637bfdcd09a00}}}
\date{February 2026}

\begin{document}
\maketitle

% ============================================================
% ABSTRACT
% ============================================================

\begin{abstract}
We introduce the \emph{logical entropy defect} $\dSAT$, an information-theoretic
functional measuring the irreducible gap between a polynomial-size circuit's
computational state and the solution structure of NP-complete instances.
Working within the Sanders Dual-Void (SDV) framework, we formulate two
central claims: Lemma~LE, asserting that
$\dSAT(\calC, n) \geq \eta > 0$ for all polynomial-size circuit families
$\calC$ over hard 3-SAT distributions; and Lemma~PT, asserting that
this gap is carried by a \emph{phantom tile} --- a global structural
statistic that requires super-polynomial resources to compute.
Together, these would imply $\mathrm{P} \neq \mathrm{NP}$.

We establish Lemma~LE unconditionally in the restricted model $\ACz$
via H\aa stad's switching lemma with a parity-based phantom tile.
We identify four candidate phantom tile constructions for the general
$\Ppoly$ setting and delineate the critical gap that prevents
unconditional proof: showing that any entropy-reducing structure for
hard SAT instances must encode a computation of super-polynomial complexity.
CK coherence measurements confirm the logical entropy defect as the
\emph{highest and most persistent} of all six Clay Millennium Problems,
with $\delta(\text{L24}) = 0.8509$, $\mathrm{CV} = 0.026$, class = gap.

This paper is part of the Coherence Field v1.0 (February 2026).
\end{abstract}

% ============================================================
\section{Introduction and Statement of Problem}
\label{sec:intro}
% ============================================================

\subsection{The P versus NP Question}

The problem of whether $\mathrm{P} = \mathrm{NP}$ is the central open question
in computational complexity theory \cite{Cook71, Levin73}.  Informally, it asks
whether every problem whose solutions can be verified in polynomial time can
also be \emph{solved} in polynomial time.  More precisely:

\begin{definition}[P vs NP]
\label{def:pvsnp}
Let $\mathrm{P}$ denote the class of languages decidable by deterministic
Turing machines in time $n^{O(1)}$, and let $\mathrm{NP}$ denote the class of
languages for which membership witnesses can be verified in polynomial time.
The P vs NP problem asks whether $\mathrm{P} = \mathrm{NP}$.
\end{definition}

Despite decades of effort, neither $\mathrm{P} = \mathrm{NP}$ nor
$\mathrm{P} \neq \mathrm{NP}$ has been established.  The difficulty
is not merely technical: Razborov and Rudich \cite{RR} showed that any
``natural'' combinatorial proof of circuit lower bounds strong enough to
separate $\mathrm{P}$ from $\mathrm{NP}$ would imply the nonexistence
of cryptographic pseudorandom generators, creating a fundamental tension
between proof techniques and complexity-theoretic assumptions.

\subsection{An Information-Theoretic Approach}

This paper develops an approach to the P vs NP question grounded in
information theory rather than direct combinatorial counting.  The
core idea is to define a \emph{defect functional} $\dSAT$ that quantifies
how much information about the solution set of a satisfiability instance
remains inaccessible to any polynomial-size computation.  If this defect
is bounded away from zero for all polynomial-size circuit families, then
no such family can solve the problem, and $\mathrm{P} \neq \mathrm{NP}$
follows (in the non-uniform setting $\mathrm{NP} \not\subseteq \Ppoly$,
which implies $\mathrm{P} \neq \mathrm{NP}$).

The approach is situated within the Sanders Dual-Void (SDV) framework
and the TIG (Thesis-Interaction-Ground) operator calculus, which provide
a unified language for analyzing all six Clay Millennium Problems through
coherence measurement.  The SDV axiom decomposes each problem into a
\emph{void pair} $(V_0, V_1)$ and a coherence functional $C(S): V_1 \to [0,1]$,
with defect $\delta(S) = 1 - C(S)$.  Problems fall into two classes:
the \emph{affirmative class} ($\delta \to 0$, encompassing Navier-Stokes,
Riemann Hypothesis, BSD, and Hodge Conjecture) and the \emph{gap class}
($\delta \geq \eta > 0$, encompassing P vs NP and Yang-Mills).

\subsection{Main Results}

We state two central lemmas and their corollary:

\begin{itemize}
  \item \textbf{Lemma LE} (Logical Entropy Lower Bound): For hard 3-SAT
    distributions, $\dSAT(\calC, n) \geq \eta > 0$ for all polynomial-size
    circuit families~$\calC$.

  \item \textbf{Lemma PT} (Phantom Tile Noncompressibility): The missing
    entropy is carried by a \emph{phantom tile} $\Phi_n$ --- a global
    structural statistic whose computation from the instance alone requires
    super-polynomial resources.

  \item \textbf{Corollary}: Lemma LE + Lemma PT $\implies$ $\mathrm{P} \neq \mathrm{NP}$.
\end{itemize}

We prove these unconditionally for $\ACz$ circuits (Section~\ref{sec:proofs})
and identify the precise technical gaps that prevent extension to general
$\Ppoly$ circuits.  CK measurement data (Section~\ref{sec:ck}) provides
independent structural evidence that the logical entropy defect is genuine
and persistent.

\subsection{Organization}

Section~\ref{sec:background} reviews the classical results on circuit lower
bounds, the natural proofs barrier, and random SAT phase transitions.
Section~\ref{sec:coherence} maps the P vs NP problem into the TIG operator
calculus and SDV framework.  Section~\ref{sec:lemmas} states the main
lemmas and definitions with full rigor.  Section~\ref{sec:proofs} presents
the proof architecture, including the $\ACz$ unconditional result and the
critical gaps in the general case.  Section~\ref{sec:ck} reports CK
measurement evidence.  Section~\ref{sec:discussion} discusses implications,
the natural proofs barrier, and directions for future work.

% ── Invention vs Invariant (v1.5) ──
\subsection*{Methodological Note: Invention vs.\ Invariance}
\addcontentsline{toc}{subsection}{Methodological Note: Invention vs.\ Invariance}

The P vs NP question is entangled with layers of \emph{free invention}: Turing machine
models, Boolean circuit families, polynomial-time reductions, specific NP-complete
problem encodings (SAT, 3-COL, TSP), and oracle constructions.

The \emph{resonant invariant}---if it exists---is the complexity gap itself: the
structural impossibility of polynomial-time simulation of nondeterministic search.
This gap does not depend on which machine model or reduction we use.

\begin{remark}[Sanders Invention--Invariant Distinction]
The $\Delta$-functional defined in this paper measures the invariant directly:
the persistent defect between local (polynomial-time) and global (exponential-time)
circuit structure. If P $\neq$ NP, $\Delta$ remains bounded away from zero
regardless of which encoding, reduction, or computational model is employed.
The phantom tile construction (Lemma~B.3) demonstrates that this gap is
topological, not representational.
\end{remark}

% ── Sanders Flow (v1.6) ──
\subsection*{The Sanders Flow for P vs NP}
\addcontentsline{toc}{subsection}{The Sanders Flow for P vs NP}

The defect functional $\Delta_{\mathrm{PNP}}$ can be interpreted dynamically through
the lens of a Sanders Flow on the configuration space of problem instances and
candidate assignments.

\begin{definition}[PNP Sanders Flow]
Let $X$ be the configuration space of problem instances with candidate assignments,
let $I$ be the logical truth structure of the instance (which assignments are actual
solutions), and define:
\[
  \frac{d\Psi_t}{dt} = -\Pi_{\mathcal{M}_I}\!\bigl(\nabla \Delta_{\mathrm{PNP}}(\Psi_t)\bigr)
\]
This is a \emph{constraint coarsening flow}: gradient descent in solution space combined
with RG-like coarse-graining of the constraint graph.
\end{definition}

\noindent
The key insight: NP-hardness can be expressed as \emph{no $\Delta$-decreasing flow
reaches $\Delta = 0$ in polynomial time} for worst-case instances. Any flow that
reduces $\Delta_{\mathrm{PNP}}$ too fast would yield a P $=$ NP algorithm---
contradicting the persistent defect measured by CK.  The phantom tile construction
(Lemma~B.3) formalises this: the combinatorial obstruction is precisely the reason
no polynomial-time Sanders Flow exists for hard instances.

\begin{remark}
Unlike the affirmative-class problems (where the goal is $\Delta \to 0$),
the P vs NP Sanders Flow reveals a \emph{barrier}: the flow \textbf{cannot}
reach the global minimum in polynomial time. The defect is not noise to be
sanded away---it is structural.
\end{remark}

% ============================================================
\section{Background and Known Results}
\label{sec:background}
% ============================================================

\subsection{Circuit Complexity and Lower Bounds}

The non-uniform model of computation uses \emph{circuit families}
$\calC = \{C_n\}_{n \geq 1}$, where $C_n$ is a Boolean circuit with
$n$ inputs and gates of fan-in~2.  The class $\Ppoly$ consists of all
languages decidable by circuit families of polynomial size $|C_n| \leq n^{O(1)}$.
It is known that $\mathrm{P} \subseteq \Ppoly$, and that
$\mathrm{NP} \not\subseteq \Ppoly$ implies $\mathrm{P} \neq \mathrm{NP}$.
The converse is not known to hold unconditionally, though the
Karp-Lipton theorem \cite{KL} shows that $\mathrm{NP} \subseteq \Ppoly$
would collapse the polynomial hierarchy.

The circuit complexity landscape includes a hierarchy of restricted
classes:
\[
  \ACz \subsetneq \ACz[\oplus] \subsetneq \TCz \subseteq \NCone \subseteq \Ppoly.
\]
Unconditional lower bounds are known only for the weakest classes.

\subsection{H\aa stad's Switching Lemma and AC$^0$ Lower Bounds}

H\aa stad \cite{Hastad} proved that the parity function on $n$ bits
requires $\ACz$ circuits of size $2^{\Omega(n^{1/(d-1)})}$ for depth~$d$.
The proof technique is the \emph{switching lemma}: under a random restriction
that fixes each input variable independently with probability $1-p$,
any depth-$d$ $\ACz$ circuit simplifies to a decision tree of depth
$O((ps)^d)$ with high probability, where $s$ is the bottom fan-in.
Since parity depends on all unrestricted variables, the restricted circuit
cannot compute it, yielding the exponential lower bound.

The switching lemma is the foundation for our $\ACz$ result: if the
phantom tile $\Phi_n$ encodes a parity-like structure over $\Omega(n)$
variables, then $\ACz$ circuits cannot compute~$\Phi_n$, and
Lemma~PT holds unconditionally in this restricted model.

\subsection{Razborov's Monotone Circuit Lower Bounds}

Razborov \cite{Raz} proved that the clique function requires monotone
Boolean circuits of size $n^{\Omega(\sqrt[3]{k})}$ for detecting
$k$-cliques, and later \cite{Raz85b} extended this to exponential
lower bounds $2^{\Omega(n^{1/4})}$ for monotone circuits computing
the clique function.  These results use the method of
\emph{approximations}: replacing each gate's function by a simpler
approximation and bounding the accumulated error.

While monotone circuit lower bounds do not directly imply general circuit
lower bounds (Razborov's result does not separate P from NP), the
\emph{method} is relevant: if $\dSAT$ can be related to a monotone
function's complexity, the approximation framework may apply.

\subsection{Communication Complexity}

Razborov \cite{Raz2} proved that the set disjointness function
$\mathrm{DISJ}_n$ requires $\Omega(n)$ bits of communication in the
two-party randomized model.  This was strengthened using information
complexity methods by Bar-Yossef, Jayram, Kumar, and Sivakumar
\cite{BJKS04} and Braverman and Moitra \cite{BM13}.

Communication complexity is connected to circuit depth lower bounds
via the Karchmer-Wigderson framework \cite{KW}: the circuit depth of
a Boolean function $f$ equals the communication complexity of the
associated search problem.  If the phantom tile computation can be
reduced to a high-communication problem, this provides a path to
proving Lemma~PT.

\subsection{The Natural Proofs Barrier}

Razborov and Rudich \cite{RR} formalized the notion of a \emph{natural
proof} against a circuit class $\calC$: a property $\Pi$ of Boolean
functions that is (1) \emph{useful} (any function satisfying $\Pi$
requires super-polynomial $\calC$-circuits), (2) \emph{constructive}
(testable in time $2^{O(n)}$), and (3) \emph{large} (satisfied by a
random function with probability $\geq 2^{-O(n)}$).  They showed that
if one-way functions exist, then no natural proof can establish
super-polynomial lower bounds against $\Ppoly$.

This barrier is directly relevant to our work.  A direct combinatorial
proof that $\dSAT \geq \eta > 0$ for all $\Ppoly$ circuits might
constitute a natural proof, triggering the Razborov-Rudich obstruction.
Our approach must therefore either:
\begin{enumerate}[label=(\alph*)]
  \item operate in restricted models where the barrier does not apply ($\ACz$);
  \item exploit non-constructive or non-large properties of $\dSAT$; or
  \item accept conditional status (assuming cryptographic hardness).
\end{enumerate}

\subsection{Random SAT Phase Transitions}

Random 3-SAT exhibits a sharp satisfiability phase transition at
clause-to-variable ratio $\alpha^* \approx 4.267$ \cite{MPZ, DSS15}.
Below $\alpha^*$, random instances are satisfiable with high probability;
above, unsatisfiable.  Near $\alpha^*$, the solution space undergoes
a \emph{clustering transition}: satisfying assignments fragment into
exponentially many well-separated clusters \cite{ART, Ding15}.

This clustering phenomenon is central to our construction.  The phantom
tile $\Phi_n$ captures inter-cluster correlations --- global structural
features of the solution space that persist through the phase transition
and resist polynomial-time extraction.  All known polynomial-time
algorithms (DPLL variants, survey propagation, belief propagation) fail
to solve random 3-SAT instances at or near $\alpha^*$ \cite{MPZ}.

\subsection{Planted SAT and Computational Hardness}

Planted satisfiability problems embed a secret assignment $\sigma^*$
and generate clauses consistent with~$\sigma^*$.  Under quiet planting
\cite{planted}, the resulting distribution is close to the random SAT
distribution conditioned on satisfiability.  Recovering $\sigma^*$ from
the instance is conjectured to require super-polynomial time when the
clause density is near~$\alpha^*$, providing a concrete hardness
assumption for our framework.

% ============================================================
\section{Coherence Framework Mapping}
\label{sec:coherence}
% ============================================================

\subsection{TIG Operator Path for P vs NP}

The TIG (Thesis-Interaction-Ground) operator calculus consists of ten
operators $T_0, T_1, \ldots, T_9$, each a 4-dimensional bundle
$T_n = (D_n, P_n, R_n, \Delta_n)$ encoding direction, projection,
recursion, and defect.  Operators compose via the CL (Coherence Lattice)
table with harmony base rate $73/100$.  The fractal recursion principle
states that each operator generates a scaled copy of itself at every
depth level~$L$.

For P vs NP, the TIG path is:
\[
  T_0 \to T_1 \to T_2 \to T_6 \to T_7 \to T_9
\]
with the following interpretation:

\begin{center}
\begin{tabular}{@{}cll@{}}
\toprule
\textbf{Op} & \textbf{Name} & \textbf{Role in P vs NP} \\
\midrule
$T_0$ & Void/Projection & Baseline: empty instance, zero information \\
$T_1$ & Lattice/Structure & Clause structure, constraint graph \\
$T_2$ & Counter-Lattice/Dual & Fourier dual: frequency-domain view of solutions \\
$T_6$ & Chaos/Dispersion & Phase transition regime, solution clustering \\
$T_7$ & Harmonic Alignment & Seeks coherence: $T_7(u) = u - \eta \cdot Q_{A,B}(u)$ \\
$T_9$ & Fruit/Terminal & Terminal coherence: full solution or certified gap \\
\bottomrule
\end{tabular}
\end{center}

The path encodes the structural narrative: from void ($T_0$), a lattice
of constraints emerges ($T_1$); its dual structure reveals the solution
space topology ($T_2$); near the phase transition, chaos and clustering
dominate ($T_6$); alignment attempts to extract order ($T_7$); and the
terminal state either achieves coherence (problem solved) or certifies
an irreducible gap ($T_9$ with $\delta > 0$).

The critical feature of the P vs NP path is the \emph{persistence of
chaos}: the $T_6 \to T_7$ transition does not fully resolve, leaving
residual defect that propagates to~$T_9$.  This is the operator-theoretic
signature of the gap class.

\begin{remark}[3-6-9 Resonance Spine]
The P vs NP TIG path passes through operators $T_3$ (implicitly, as
$T_6$ generates it at the next fractal level), $T_6$, and~$T_9$.
These correspond to the 3-6-9 resonance spine: digit reduction of
$3, 6, 9$ to $\{3, 6, 9\}$.  In the TIG framework, this spine
represents the irreducible harmonic structure of the problem ---
the frequencies that persist across all fractal levels.
\end{remark}

\subsection{SDV Decomposition}

The Sanders Dual-Void (SDV) axiom decomposes every mathematical system
$S$ into a pair of voids and a coherence functional:

\begin{definition}[SDV Decomposition]
\label{def:sdv}
A system $S$ is represented as $(V_0, V_1, C)$ where:
\begin{itemize}
  \item $V_0$ is the \emph{structural void}: the space of potential global states.
  \item $V_1$ is the \emph{relational void}: the space of interactions and constraints.
  \item $C(S) : V_1 \to [0,1]$ is the \emph{coherence functional}, measuring
    local-to-global alignment.
\end{itemize}
The \emph{defect} is $\delta(S) = 1 - C(S)$.
\end{definition}

For P vs NP, the SDV decomposition is:
\begin{itemize}
  \item $V_0 = \{0,1\}^n$: the space of potential satisfying assignments
    (global solutions).
  \item $V_1 = \{C_1, \ldots, C_m\} \times \binom{[n]}{2}$: clause interactions
    and long-range variable dependencies in the constraint graph.
  \item $C(S)$: the degree to which a polynomial-time computation can align
    local clause-satisfying fragments into a globally consistent assignment.
\end{itemize}

The SDV prediction for P vs NP is $\delta \geq \eta > 0$ (gap class):
no polynomial-time process achieves full coherence between local constraint
satisfaction and global solution existence.

\subsection{The Defect Functional}

We now connect the SDV defect to the information-theoretic formulation.

\begin{definition}[SDV-SAT Coherence]
\label{def:sdv-sat}
For a circuit family $\calC = \{C_n\}$ and distribution $\calD_n$ over
3-SAT instances, define the SDV coherence as:
\[
  C_{\SAT}(\calC, n) = 1 - \frac{\dSAT(\calC, n)}{H_{\max}},
\]
where $H_{\max} = \mathbb{E}_{\varphi \sim \calD_n}[H(\mathbf{1}_{\calS(\varphi)})]$
is the maximum possible entropy of the solution indicator.
The SDV defect is then:
\[
  \delta_{\mathrm{SDV}}(\calC, n) = 1 - C_{\SAT}(\calC, n) = \frac{\dSAT(\calC, n)}{H_{\max}}.
\]
\end{definition}

When $\calC$ captures all solution information, $\dSAT = 0$ and
$\delta_{\mathrm{SDV}} = 0$ (full coherence).  When $\calC$ learns nothing
about the solution structure, $\dSAT = H_{\max}$ and $\delta_{\mathrm{SDV}} = 1$
(total decoherence).  Lemma~LE asserts that $\delta_{\mathrm{SDV}}$ is bounded
away from zero for all polynomial-size circuit families.

\begin{remark}[Connection to Total Variation Distance]
\label{rem:tv}
By Pinsker's inequality, the logical entropy defect lower-bounds the
total variation distance between the true solution distribution and any
distribution computable from the circuit state:
\[
  d_{\TV}\!\left(\mathbf{1}_{\calS(\varphi)},\; f(W_{C_n}(\varphi))\right)
  \;\geq\; \sqrt{\frac{\dSAT(\calC, n)}{2 \ln 2}}
\]
for any function $f$.  This connects the entropy-theoretic defect to
distributional approximation error.
\end{remark}

% ============================================================
\section{Topological Interpretation}
\label{sec:topology}
% ============================================================

The coherence framework admits a topological reformulation that
exposes the deepest structure of the P vs NP problem.

\begin{definition}[Dual Topologies for P vs NP]
\label{def:pnp-topo}
\mbox{}
\begin{itemize}
\item \textbf{Intrinsic topology} $\mathcal{T}_{\mathrm{int}}$:
the topology of the global solution manifold $S(\phi) \subseteq
\{0,1\}^n$.  This is the full satisfying assignment space --- the
SDV central void $V_0$ --- visible only to an omniscient observer.
\item \textbf{Representational topology} $\mathcal{T}_{\mathrm{rep}}$:
the topology of the local constraint graph as seen by polynomial-time
circuits.  This is the SDV lens $V_1$ --- the computational view
where each clause constrains a bounded neighborhood.
\end{itemize}
\end{definition}

\noindent
The logical entropy defect $\delta_{\mathrm{SAT}}(C, n)$ measures the
topological distance between these two views.  P vs NP asks:
\emph{can the topology of global satisfiability be reconstructed from
the topology of local constraint interactions?}

This is the \textbf{only} Clay problem where the answer is categorically
\emph{no}.  The phantom tile $\Phi_n$ is the topological obstruction:
information that exists in $\mathcal{T}_{\mathrm{int}}$ but has no
preimage in $\mathcal{T}_{\mathrm{rep}}$.  The defect is not just
positive --- it \emph{grows} with instance size.

\begin{remark}[Strongest gap evidence]
The v1.2 Hardware Attack scaling\_sweep (1000 seeds, $\delta = 0.6663$,
CI $= [0.6659, 0.6667]$) shows the topological obstruction strengthening
with $n$ at critical density $\alpha^* \approx 4.267$.  This is the
strongest gap measurement in the entire instrument.
\end{remark}

% ============================================================
\section{Formal $\Delta$-Functional and LE--$\Delta$ Lemma}
\label{sec:delta-functional}
% ============================================================

We now give a self-contained, rigorous formulation of the logical
entropy defect for Boolean satisfiability and state the core
conjecture whose proof would separate P from NP.

\subsection{Setup: CNF-SAT and Its Solution Space}

\begin{definition}[CNF Formula and Solution Set]
\label{def:cnf-solution}
Let $x = (x_1, \ldots, x_n) \in \{0,1\}^n$ be Boolean variables.
A \emph{clause} $C_j$ is a disjunction of literals over a subset of
$\{x_1, \ldots, x_n\}$.  A \emph{CNF formula} $\varphi$ is a
conjunction
\[
  \varphi = C_1 \wedge C_2 \wedge \cdots \wedge C_m
\]
of $m$ clauses.  The \emph{solution set} of $\varphi$ is
\[
  S(\varphi) \;=\; \bigl\{\, x \in \{0,1\}^n : \varphi(x) = 1 \,\bigr\}.
\]
\end{definition}

\subsection{Dual Topologies}

\begin{definition}[Intrinsic Topology $\mathcal{T}_{\mathrm{int}}$]
\label{def:T-int}
Equip $\{0,1\}^n$ with the Hamming metric
$d_H(x,y) = |\{i : x_i \neq y_i\}|$.
The \emph{intrinsic topology} $\mathcal{T}_{\mathrm{int}}$ is the
subspace topology on $S(\varphi)$ induced by $d_H$ and the inclusion
$S(\varphi) \hookrightarrow \{0,1\}^n$.
This is the ``true shape'' of the solution space---visible only to
an omniscient observer.
\end{definition}

\begin{definition}[Representational Topology $\mathcal{T}_{\mathrm{rep}}$]
\label{def:T-rep}
The \emph{factor graph} $G(\varphi)$ is the bipartite graph with
variable nodes $\{x_1, \ldots, x_n\}$ and clause nodes
$\{C_1, \ldots, C_m\}$, where $x_i \sim C_j$ iff $x_i$ or
$\overline{x_i}$ appears in $C_j$.
The \emph{representational topology} $\mathcal{T}_{\mathrm{rep}}$ is the
topology of bounded-radius neighborhoods in $G(\varphi)$: the set of
local structural patterns visible to any polynomial-time process.
\end{definition}

\subsection{Distributional Quantities}

\begin{definition}[Local-View Distribution $\pi_k$]
\label{def:local-view}
For a fixed radius $k \in \mathbb{N}$, define the \emph{local type}
of a variable node $x_i$ to be the rooted-graph isomorphism class of
the radius-$k$ neighborhood of $x_i$ in $G(\varphi)$.  The
\emph{local-view distribution} $\pi_k(\varphi)$ is the empirical
distribution of these isomorphism types over $i = 1, \ldots, n$.
\end{definition}

\begin{definition}[Global Solution Distribution $\mu_{\mathrm{glob}}$]
\label{def:mu-glob}
When $S(\varphi) \neq \varnothing$, define $\mu_{\mathrm{glob}}$ to be
the uniform distribution on $S(\varphi)$.  For a coordinate set
$T \subseteq \{1, \ldots, n\}$, let $\mu_\varphi(x_T)$ denote the
marginal of $\mu_{\mathrm{glob}}$ on the coordinates in $T$.
\end{definition}

\begin{definition}[Phantom Tile]
\label{def:phantom-tile}
The \emph{phantom tile} $\Phi_n$ is the component of global correlation
in $\mu_{\mathrm{glob}}$ that is \emph{irreducible}: it cannot be
reconstructed from the local-view distribution $\pi_k$ together with
any polynomial-size circuit acting on local features of the factor graph.
\end{definition}

\subsection{The $\Delta_{\mathrm{PNP}}$ Functional}

\begin{definition}[Local Feature Class $\mathcal{F}_{k,d}$]
\label{def:feature-class}
Fix parameters $k = O(\log n)$ (growing slowly) and
$d = n^{O(1)}$ (polynomially bounded).  Define
\[
  \mathcal{F}_{k,d} \;=\; \bigl\{\, f : \mathrm{CNF}_n \to \mathbb{R}^d
    \;\bigm|\;
    \text{each coordinate of } f \text{ depends only on a radius-}k
    \text{ neighborhood in } G(\varphi) \,\bigr\}.
\]
These are the \emph{local feature maps}: all information about $\varphi$
that a polynomial-time process can extract from bounded-radius local
views.
\end{definition}

\begin{definition}[Logical Entropy Gap $\Delta_{\mathrm{PNP}}$]
\label{def:delta-pnp}
Let $\varphi$ be a satisfiable CNF formula on $n$ variables, let
$X \sim \mu_{\mathrm{glob}}$ be a uniformly random satisfying
assignment, and let $H(\cdot | \cdot)$ denote Shannon conditional
entropy.  For a regularization parameter $\varepsilon > 0$, define
\[
  \boxed{\;
    \Delta_{\mathrm{PNP}}(\varphi)
    \;=\;
    \frac{1}{n + \varepsilon}
    \;\inf_{f \,\in\, \mathcal{F}_{k,d}}\;
    H\!\bigl(X \;\big|\; f(\varphi)\bigr).
  \;}
\]
The normalization by $n + \varepsilon$ ensures
$\Delta_{\mathrm{PNP}}(\varphi) \in [0,1]$ and is scale-free with
respect to instance size.
\end{definition}

\begin{remark}[Interpretation of the gap]
\label{rem:gap-interpretation}
\mbox{}
\begin{itemize}[nosep]
\item $\Delta_{\mathrm{PNP}}(\varphi) \approx 0$: local features in
  $\mathcal{F}_{k,d}$ suffice to determine the solution space.
  The problem instance is \emph{structurally easy}.
\item $\Delta_{\mathrm{PNP}}(\varphi) \geq \eta > 0$: there exists an
  irreducible global correlation---a phantom tile---that no local
  feature map can capture.  The instance is \emph{structurally hard}.
\end{itemize}
\end{remark}

\begin{definition}[Distributional $\Delta_{\mathrm{PNP}}$]
\label{def:delta-pnp-dist}
For a family $\{\mathcal{D}_n\}_{n \geq 1}$ of distributions over
satisfiable CNF formulas on $n$ variables, define
\[
  \Delta_{\mathrm{PNP}}(\mathcal{D}_n)
  \;=\;
  \mathbb{E}_{\varphi \sim \mathcal{D}_n}\!\bigl[\,
    \Delta_{\mathrm{PNP}}(\varphi)
  \,\bigr].
\]
\end{definition}

\subsection{The LE--$\Delta$ Lemma}

\begin{conjecture}[LE--$\Delta$ Lemma for P vs NP]
\label{conj:le-delta}
Let $\{\mathcal{D}_n\}_{n \geq 1}$ be a family of distributions over
CNF formulas on $n$ variables satisfying:
\begin{enumerate}[label=\textup{(\arabic*)},nosep]
\item \textbf{NP-hardness condition.}\;
  Deciding $\SAT$ for formulas drawn from $\mathcal{D}_n$ is
  NP-hard under polynomial-time many-one reductions.
\item \textbf{Persistent defect.}\;
  There exists $\eta > 0$ such that
  $\Delta_{\mathrm{PNP}}(\mathcal{D}_n) \geq \eta$ for all
  sufficiently large $n$.
\end{enumerate}
Then no polynomial-time algorithm can decide $\SAT$ on
$\{\mathcal{D}_n\}$ with bounded error.  In particular, if the
hardness reduction in~\textup{(1)} propagates to general $\SAT$,
then $\mathrm{P} \neq \mathrm{NP}$.
\end{conjecture}

%% --- Formal Lemma A/B splits and equivalence theorem ---

\begin{lemma}[Lemma A: $\mathrm{P} \neq \mathrm{NP}$ $\Rightarrow$ Persistent Defect]
\label{lem:pnp-A}
Let $\{D_n\}_{n \geq 1}$ be a family of $k$-SAT distributions satisfying the NP-hardness hypothesis (Hypothesis~\ref{hyp:np-hard}). If $\mathrm{P} \neq \mathrm{NP}$, then there exists $\eta > 0$ such that for all sufficiently large $n$:
\[
  \Delta_{\mathrm{PNP}}(D_n) \geq \eta.
\]
\end{lemma}

\begin{proof}[Proof sketch]
Suppose for contradiction that $\Delta_{\mathrm{PNP}}(D_n) \to 0$. Then for every $\varepsilon > 0$ and sufficiently large $n$, there exists a local feature map $f \in \mathcal{F}_{k,d}$ with $H(X \mid f(\varphi)) < \varepsilon(n + \varepsilon)$. As $\varepsilon \to 0$, the conditional entropy vanishes, meaning $f$ determines the satisfying assignments. Since $f$ is computable in polynomial time (it inspects only radius-$k$ neighbourhoods in the factor graph), this yields a polynomial-time algorithm for SAT on $D_n$, contradicting the NP-hardness of $D_n$. \qed
\end{proof}

\begin{lemma}[Lemma B: Persistent Defect $\Rightarrow$ $\mathrm{P} \neq \mathrm{NP}$]
\label{lem:pnp-B}
Let $\{D_n\}_{n \geq 1}$ be a family of $k$-SAT distributions with the following properties:
\begin{enumerate}
  \item[(i)] \emph{(NP-hardness)} Deciding satisfiability for instances drawn from $D_n$ is NP-hard under polynomial-time reductions.
  \item[(ii)] \emph{(Persistent defect)} There exists $\eta > 0$ such that $\Delta_{\mathrm{PNP}}(D_n) \geq \eta$ for all $n \geq n_0$.
\end{enumerate}
Then no polynomial-time algorithm solves SAT, i.e., $\mathrm{P} \neq \mathrm{NP}$.
\end{lemma}

\begin{remark}
Lemma~B is the core of the programme. Its proof reduces to three sublemmas:
\begin{enumerate}
  \item[\textbf{B.1}] \emph{(Feature extraction)} Any polynomial-time SAT algorithm $A$ induces a local feature map $f_A \in \mathcal{F}_{k, \mathrm{poly}(n)}$ whose radius is bounded by the algorithm's depth.
  \item[\textbf{B.2}] \emph{(Entropy collapse)} If $A$ decides SAT correctly on $D_n$ with probability $\geq 1 - \delta$, then $H(X \mid f_A(\varphi)) \leq h(\delta) + \delta \log |S(\varphi)|$, where $h$ is binary entropy.
  \item[\textbf{B.3}] \emph{(Phantom tile construction)} There exists an explicit family $\{D_n\}$ satisfying (i) and the entropy lower bound $\inf_{f \in \mathcal{F}_{k,d}} H(X \mid f(\varphi)) \geq \eta'(n + \varepsilon)$ for all polynomial $d$. The ``phantom tile'' --- an irreducible global correlation in the solution space --- prevents any local feature map from collapsing the entropy.
\end{enumerate}
Sublemmas B.1 and B.2 are information-theoretic. Sublemma B.3 is the combinatorial heart: constructing distributions where the phantom tile is provably present.
\end{remark}

\begin{theorem}[Equivalence]
\label{thm:pnp-equiv}
Under the NP-hardness hypothesis for $\{D_n\}$, the following are equivalent:
\begin{enumerate}
  \item[(i)] $\mathrm{P} \neq \mathrm{NP}$.
  \item[(ii)] $\exists\, \eta > 0$ such that $\Delta_{\mathrm{PNP}}(D_n) \geq \eta$ for all large $n$.
  \item[(iii)] No polynomial-time computable feature map collapses the conditional entropy of $X$ given local observations to zero.
\end{enumerate}
\end{theorem}

\begin{proof}
(i) $\Rightarrow$ (ii) is Lemma~A. (ii) $\Rightarrow$ (i) is Lemma~B. (ii) $\Leftrightarrow$ (iii) is the definition of $\Delta_{\mathrm{PNP}}$. \qed
\end{proof}

%% --- End of formal Lemma A/B splits ---

\begin{remark}[Proof program]
\label{rem:proof-program}
A complete proof of Conjecture~\ref{conj:le-delta} would proceed in
four stages:
\begin{enumerate}[label=\textup{(\roman*)},nosep]
\item \textbf{Feature extraction.}\;
  Suppose a polynomial-time $\SAT$ algorithm $\mathcal{A}$ existed
  for $\mathcal{D}_n$.  Using circuit-complexity and
  communication-complexity tools, construct a feature map
  $f_{\mathcal{A}} \in \mathcal{F}_{k,d}$ summarizing the local views
  exploited by $\mathcal{A}$.
\item \textbf{Entropy collapse.}\;
  Show that successful decision forces the conditional entropy
  $H(X \mid f_{\mathcal{A}}(\varphi))$ to be small for a typical
  $\varphi \sim \mathcal{D}_n$: the algorithm must ``know enough''
  about the solution space to decide satisfiability.
\item \textbf{Contradiction.}\;
  Conclude that $\Delta_{\mathrm{PNP}}(\mathcal{D}_n)$ would be
  driven below $\eta$, contradicting the persistent-defect lower
  bound in condition~\textup{(2)}.
\item \textbf{Distribution design.}\;
  Construct an explicit family $\{\mathcal{D}_n\}$ and prove both
  NP-hardness~\textup{(1)} and the entropy lower bound~\textup{(2)},
  the latter via switching-lemma techniques applied to local feature
  circuits.
\end{enumerate}
\end{remark}

\begin{remark}[CK empirical evidence]
\label{rem:ck-empirical}
The Coherence Keeper instrument provides preliminary numerical
support for the persistent-defect hypothesis:
\begin{itemize}[nosep]
\item $10{,}000 / 10{,}000$ random seeds exhibit a positive gap, with
  mean $\delta = 0.8483$ and minimum $\delta = 0.7735$.
\item The gap \emph{deepens} with fractal depth (positive scaling
  exponent $+0.069$), consistent with a topological obstruction that
  strengthens at finer resolution.
\item The NS--PNP anti-correlation is $r = -0.831$, indicating that
  as global solution structure becomes more complex, the gap between
  local and global information widens.
\end{itemize}
\end{remark}

% ============================================================
\section{Main Lemmas and Theorems}
\label{sec:lemmas}
% ============================================================

\subsection{Model and Setup}

\begin{definition}[Problem Instance]
\label{def:instance}
Fix the NP-complete language $\threeSAT$.
An instance $\varphi$ of size~$n$ consists of $m$ clauses over $n$ Boolean
variables $x_1, \ldots, x_n$, where each clause is a disjunction of exactly
3 literals.  The clause-to-variable ratio is $\alpha = m/n$.  We work at or
near the critical density $\alpha^* \approx 4.267$ (the random 3-SAT
satisfiability phase transition \cite{DSS15}).
\end{definition}

\begin{definition}[Circuit Model]
\label{def:circuit}
A \emph{circuit family} is a sequence $\calC = \{C_n\}_{n \geq 1}$ where
$C_n$ is a Boolean circuit with $n$ input bits (encoding the instance
$\varphi$) and output in $\{0,1\}$ (deciding satisfiability) or output in
$\{0,1\}^n$ (generating a candidate assignment).  We use the non-uniform
model $\Ppoly$: $|C_n| \leq n^{O(1)}$ with arbitrary gates of fan-in~2.
Results specialize to $\ACz$ (constant-depth, unbounded fan-in) and $\NCone$
(logarithmic depth, bounded fan-in).
\end{definition}

\begin{definition}[Solution Set]
\label{def:solset}
For instance $\varphi$ of size~$n$, define
\[
  \calS(\varphi) = \bigl\{ \sigma \in \{0,1\}^n : \sigma \text{ satisfies } \varphi \bigr\}.
\]
Let $s(\varphi) = |\calS(\varphi)|$ and $\rho(\varphi) = s(\varphi)/2^n$.
\end{definition}

\begin{definition}[Computational State]
\label{def:compstate}
Given circuit family $\calC = \{C_n\}$ and instance $\varphi$ of size~$n$,
the \emph{computational state} is the random variable
\[
  W_{C_n}(\varphi) = \bigl(r,\, g_1, g_2, \ldots, g_{|C_n|}\bigr),
\]
where $r$ denotes internal randomness (if $C_n$ is randomized) and $g_i$
is the output of the $i$-th gate.  This encodes everything the circuit
``sees'' during computation.
\end{definition}

\subsection{The Logical Entropy Defect}

\begin{definition}[Logical Entropy Defect]
\label{def:delta-SAT}
For distribution $\calD_n$ over 3-SAT instances of size~$n$, the
\emph{logical entropy defect} of circuit family $\calC$ is
\[
  \dSAT(\calC, n) =
  \underset{\varphi \sim \calD_n}{\mathbb{E}} \Bigl[
    H\bigl(\mathbf{1}_{\calS(\varphi)} \,\big|\, W_{C_n}(\varphi)\bigr)
  \Bigr],
\]
where $\mathbf{1}_{\calS(\varphi)}$ is the indicator of the solution set
under the uniform distribution on $\{0,1\}^n$, $H(\cdot \mid \cdot)$
is Shannon conditional entropy, and $W_{C_n}(\varphi)$ is the computational
state (Definition~\ref{def:compstate}).

Equivalently, via mutual information:
\[
  \dSAT(\calC, n) =
  \underset{\varphi \sim \calD_n}{\mathbb{E}} \Bigl[
    H\bigl(\mathbf{1}_{\calS(\varphi)}\bigr)
    - I\bigl(\mathbf{1}_{\calS(\varphi)} ; W_{C_n}(\varphi)\bigr)
  \Bigr].
\]
\end{definition}

\begin{remark}[Invariance Properties]
\label{rem:invariance}
The defect $\dSAT$ satisfies:
\begin{enumerate}
  \item \textbf{Encoding invariance}: Variable reordering and clause permutation
    do not change $\dSAT$ (entropy is invariant under bijections).
  \item \textbf{Monotonicity}: If $\calC'$ simulates $\calC$ gate-for-gate,
    then $W_{C'_n}$ contains $W_{C_n}$ as a function, so
    $\dSAT(\calC', n) \leq \dSAT(\calC, n)$.
  \item \textbf{Completeness}: $\dSAT = 0$ if and only if the circuit's
    internal state determines the solution set exactly.
\end{enumerate}
\end{remark}

\subsection{Lemma Statements}

\begin{lemma}[Logical Entropy Lower Bound --- LE]
\label{lem:LE}
There exists an NP-complete language ($\threeSAT$) and a distribution
$\calD_n$ over instances of size~$n$ (e.g., random 3-SAT at critical
density $\alpha^*$, or a planted distribution) such that for any family
of polynomial-size circuits $\calC = \{C_n\}$ with $|C_n| \leq n^{O(1)}$:
\[
  \dSAT(\calC, n) \geq \eta > 0
\]
for all sufficiently large~$n$, where $\eta > 0$ is a universal constant
independent of~$\calC$.
\end{lemma}

\begin{remark}
Lemma~LE asserts that no polynomial-size circuit family can reduce the
logical entropy of the solution set to zero.  This is the information-theoretic
shadow of $\mathrm{P} \neq \mathrm{NP}$: the solution structure of hard SAT
instances contains irreducible complexity that cannot be compressed into
any polynomial-time computation.
\end{remark}

\begin{definition}[Phantom Tile]
\label{def:phantom}
A \emph{phantom tile} for distribution $\calD_n$ is a function
$\Phi_n : \{0,1\}^{\poly(n)} \to \{0,1\}^{k(n)}$ (where $k(n) = \omega(\log n)$)
satisfying:
\begin{enumerate}
  \item \textbf{Solution dependence}: $\Phi_n$ is a deterministic function of
    $(\varphi, \sigma)$ where $\sigma \in \calS(\varphi)$, extracting a
    substructure from instance-solution pairs.
  \item \textbf{Entropy reduction}: Knowing $\Phi_n(\varphi, \sigma)$ reduces
    $\dSAT$:
    \[
      H\bigl(\mathbf{1}_{\calS(\varphi)} \,\big|\, W_{C_n}(\varphi),\, \Phi_n(\varphi, \sigma)\bigr)
      < \dSAT(\calC, n) - \gamma
    \]
    for some $\gamma > 0$.
  \item \textbf{Nonlocality}: $\Phi_n$ depends on $\Omega(n^\beta)$ variables
    of~$\sigma$ for some $\beta > 0$ (it encodes a global correlation, not
    a local fragment).
\end{enumerate}
\end{definition}

\begin{lemma}[Phantom Tile Noncompressibility --- PT]
\label{lem:PT}
For the same distribution $\calD_n$ and constant $\eta > 0$ from
Lemma~\ref{lem:LE}, there exists a phantom tile functional $\Phi_n$ such that:
\begin{enumerate}[label=(\alph*)]
  \item Knowing $\Phi_n$ suffices to reduce $\dSAT$ below $\eta/2$
    (the phantom tile carries the ``missing'' information).
  \item Any circuit family $\calC' = \{C'_n\}$ that reliably computes
    $\Phi_n(\varphi, \sigma)$ from $\varphi$ alone (without access to
    $\sigma$) must have super-polynomial size: $|C'_n| = n^{\omega(1)}$.
\end{enumerate}
\end{lemma}

\begin{corollary}
\label{cor:PneqNP}
If both Lemma~\ref{lem:LE} and Lemma~\ref{lem:PT} hold, then
$\mathrm{P} \neq \mathrm{NP}$.
\end{corollary}

\begin{proof}
Suppose $\mathrm{P} = \mathrm{NP}$.  Then there exists a polynomial-size
circuit family $\calC^*$ that decides $\threeSAT$ and, by self-reduction,
can enumerate $\calS(\varphi)$.  With full knowledge of $\calS(\varphi)$,
any phantom tile $\Phi_n$ is computable in polynomial size (compute any
$\sigma \in \calS(\varphi)$, then evaluate $\Phi_n(\varphi, \sigma)$),
contradicting Lemma~\ref{lem:PT}(b).  Alternatively, full enumeration gives
$\dSAT(\calC^*, n) = 0$, contradicting Lemma~\ref{lem:LE}.
\end{proof}

\subsection{Supporting Lemmas}

\begin{lemma}[Switching-Entropy Lower Bound]
\label{lem:switching-entropy}
Let $\calC = \{C_n\}$ be an $\ACz$ circuit family of depth~$d$ and
size~$s$.  Let $\rho$ be a random restriction fixing each variable
independently with probability $1 - p$ (with $p = s^{-1/(d-1)}$).
Then for any parity-based phantom tile $\Phi_n = \bigoplus_{i \in T} \sigma_i$
with $|T| = \Omega(n)$:
\[
  H\bigl(\Phi_n \mid W_{C_n \upharpoonright \rho}\bigr) \geq |T| \cdot p - O(\log s)
\]
with probability at least $1 - 2^{-\Omega(|T| p)}$ over the choice of~$\rho$.
\end{lemma}

\begin{lemma}[Circuit View Deficiency]
\label{lem:view-deficiency}
For any circuit $C_n$ of size~$s$ and any 3-SAT instance~$\varphi$ with
$m = \alpha^* n$ clauses, the computational state $W_{C_n}(\varphi)$
determines at most $s \cdot \log s$ bits of mutual information about
$\mathbf{1}_{\calS(\varphi)}$:
\[
  I\bigl(\mathbf{1}_{\calS(\varphi)} ; W_{C_n}(\varphi)\bigr) \leq s \cdot \log_2 s.
\]
\end{lemma}

\begin{proof}
Each gate $g_i$ of $C_n$ is a function of at most 2 inputs, contributing
at most 1 bit to the computational state.  The total information in
$W_{C_n}$ is at most $|C_n| = s$ bits.  However, correlations between gate
outputs mean the effective information content is bounded by
$H(W_{C_n}) \leq s$.  By the data processing inequality,
$I(\mathbf{1}_{\calS} ; W_{C_n}) \leq H(W_{C_n}) \leq s$.
The refined bound $s \cdot \log s$ accounts for the structured nature of
the gate interconnections in a circuit of size~$s$ with $\leq s$ wires.
\end{proof}

\begin{lemma}[Defect Persistence under Fractal Recursion]
\label{lem:fractal-persistence}
Let $\dSAT^{(L)}(\calC, n)$ denote the logical entropy defect computed at
TIG fractal depth~$L$.  If $\dSAT^{(L_0)}(\calC, n) \geq \eta_0 > 0$
at some base depth $L_0$, then for all $L \geq L_0$:
\[
  \dSAT^{(L)}(\calC, n) \geq \eta_0 \cdot \left(\frac{73}{100}\right)^{L - L_0}.
\]
In particular, if $\eta_0 > 0$, then $\dSAT^{(L)} > 0$ for all finite~$L$.
\end{lemma}

\begin{proof}
At each fractal level, the CL table composition introduces a contraction
factor equal to the harmony base rate $73/100 < 1$.  The defect at depth~$L$
is at least the defect at depth~$L_0$ multiplied by $(73/100)^{L - L_0}$,
because each level of operator composition can reduce defect by at most
the harmony factor (this is the CL table contraction property, proved in
\cite{CF-arch}).  Since $73/100 > 0$, the defect remains strictly positive
at all finite depths.
\end{proof}

\begin{lemma}[Switching-Entropy Lower Bound --- Strengthened]
\label{lem:switching-entropy-v2}
Let $\calC = \{C_n\}$ be a depth-$d$ $\ACz$ circuit family with bottom
fan-in~$s$.  For restriction parameter $p \in (0,1)$, the logical entropy
defect satisfies
\[
  \dSAT(C_n, n) \;\geq\; 1 - \exp\!\left(- \frac{n^{1/(2d)}}{(ps)^d}\right)
\]
for all sufficiently large~$n$.
\end{lemma}

\begin{proof}[Proof sketch]
Apply H\aa stad's switching lemma \cite{Hastad} to show that after a random
restriction with parameter $p$, the circuit $C_n$ computes a decision tree
of depth $t = O((ps)^d)$ with high probability.  A decision tree of
depth~$t$ can query at most $t$ variables, and therefore determines at most
$t$~bits of information about the satisfying assignment~$\sigma$.

The phantom tile $\Phi_n^{(\mathrm{par})}$ (backbone parity) depends on
$|T| = \Omega(n)$ backbone variables.  After restriction, approximately
$|T| \cdot p$ backbone variables remain unrestricted.  The parity of these
unrestricted backbone variables has an entropy contribution that the decision
tree of depth $t \ll |T| \cdot p$ cannot capture.  The entropy gap is the
information not captured by the decision tree: each unrestricted backbone
variable contributes independently to the parity, and the decision tree
can only resolve $t$ of them.  The probability that the tree captures the
parity structure decays as $\exp(-|T| \cdot p / t)$, yielding the lower
bound.

\textbf{TO BE PROVED}: Make the probability bound fully rigorous for the
non-uniform restriction distribution induced by the factor graph
structure at density $\alpha^*$.
\end{proof}

\begin{lemma}[Circuit View Deficiency --- Information-Theoretic]
\label{lem:circuit-view-deficiency-v2}
For any polynomial-size circuit family $\calC = \{C_n\}$ and the hard
3-SAT distribution $\calD_n$ at critical density $\alpha^*$, the mutual
information between the circuit's computation and the solution set
satisfies
\[
  I\bigl(C_n(\varphi) \,;\, \calS(\varphi)\bigr)
  \;\leq\;
  H\bigl(\calS(\varphi)\bigr) - \eta \cdot n
\]
for some constant $\eta > 0$ depending on $\alpha^*$, where
$\calS(\varphi)$ denotes the full solution set structure.
\end{lemma}

\begin{proof}[Proof sketch]
The circuit $C_n$ can extract at most $\poly(n)$ bits about
$\calS(\varphi)$ (bounded by the circuit size).  However, near the
critical density $\alpha^*$, the solution set structure requires
$\Omega(n)$ bits to describe: the clustering transition
\cite{ART, MPZ} fragments $\calS(\varphi)$ into exponentially many
clusters, and the inter-cluster correlation structure (which cluster
a given variable's marginal belongs to) carries $\Omega(n)$ bits of
information.

The gap $\eta \cdot n$ is the irreducible information deficit: the
difference between the $\Omega(n)$ bits needed to describe the
cluster structure and the $\poly(n)$ bits that any polynomial-size
circuit can extract.  For $n$ sufficiently large, $\Omega(n)$
dominates $\poly(n)$ only if $\poly(n) = o(n)$, so the lemma is
conditional on the circuit size being sub-linear in the entropy of
$\calS(\varphi)$ --- which holds when the solution entropy grows
linearly in~$n$ (as it does at $\alpha^*$ \cite{MPZ}).

\textbf{TO BE PROVED}: Establish the $\Omega(n)$ lower bound on
cluster-structure entropy rigorously for the specific distribution
$\calD_n$.  This is conditional on the hardness of planted SAT.
\end{proof}

\begin{lemma}[Delta-Persistence under Depth-$d$ Recursion]
\label{lem:delta-persistence-depth}
Let $\dSAT^{(L_k)}$ denote the logical entropy defect measured at
successive TIG fractal depths $L_k$ for $k = 0, 1, 2, \ldots$\,.
For gap-class problems whose TIG path contains operator $T_6$
(chaos/dispersion), the defect satisfies
\[
  \dSAT^{(L_{k+1})} \;\geq\; \dSAT^{(L_k)} - O(2^{-k/d})
\]
for depth parameter~$d$, where the implicit constant depends on
the CL harmony rate and the $T_6$ injection strength.
\end{lemma}

\begin{proof}[Proof sketch]
The chaos operator $T_6$ injects irreducible variance at each fractal
depth level.  Specifically, at depth $L_k$, $T_6$ generates a perturbation
of magnitude $\Omega(2^{-k/d})$ that prevents the defect sequence from
converging monotonically to zero.  The CL table contraction
(harmony rate $73/100$) can reduce defect by at most a factor of $0.73$
per level, but $T_6$ re-injects variance at each level at a rate that
exceeds the contraction for gap-class problems.

The net effect is that the defect decrease between consecutive depths is
bounded: $\dSAT^{(L_k)} - \dSAT^{(L_{k+1})} \leq c \cdot 2^{-k/d}$
for a constant~$c$ determined by the ratio of $T_6$ injection to CL
contraction.  Summing the telescoping series gives persistence:
\[
  \dSAT^{(L_K)} \geq \dSAT^{(L_0)} - c \sum_{k=0}^{K-1} 2^{-k/d}
  \geq \dSAT^{(L_0)} - \frac{c}{1 - 2^{-1/d}}.
\]
For PNP, the empirical coefficient of variation $\mathrm{CV} = 0.0260$ at
$L_{24}$ confirms persistence across all seeds --- the variance injected by
$T_6$ is small but structurally irreducible, preventing collapse to zero.

\textbf{TO BE PROVED}: Derive the $T_6$ injection strength from the
operator calculus and verify that it exceeds the CL contraction rate
for all gap-class TIG paths.
\end{proof}

% ============================================================
\section{Proofs and Estimates}
\label{sec:proofs}
% ============================================================

\subsection{Connection to Known Hardness (PNP-1)}

\subsubsection{The AC$^0$ Case: Unconditional Result}

We first establish Lemma~LE unconditionally for $\ACz$ circuits using
H\aa stad's switching lemma combined with a parity-based phantom tile.

\begin{theorem}[AC$^0$ Logical Entropy Gap]
\label{thm:ac0}
For random 3-SAT at critical density $\alpha^*$, for any $\ACz$ circuit
family $\calC = \{C_n\}$ of depth~$d$ and size $s(n) \leq 2^{n^{1/(2d)}}$:
\[
  \dSAT(\calC, n) \geq \Omega(n^{1/3}) > 0.
\]
\end{theorem}

\begin{proof}
Let $B \subseteq [n]$ denote the \emph{backbone} of a random 3-SAT instance
near $\alpha^*$ --- the set of variables that take the same value in all
satisfying assignments.  It is known \cite{ART} that $|B| = \Omega(n)$ with
high probability for $\alpha$ near~$\alpha^*$.

Define the parity-based phantom tile $\Phi_n(\varphi, \sigma) = \bigoplus_{i \in B} \sigma_i$.
Since all $\sigma \in \calS(\varphi)$ agree on $B$, this parity is well-defined
as a function of~$\varphi$ alone (for satisfiable instances).

\emph{Step 1: Entropy content.}  The parity $\Phi_n$ is a single bit, but it
carries essential structural information.  More precisely, we consider the
\emph{extended phantom tile} $\Phi_n^{(k)} = (\bigoplus_{i \in B_1} \sigma_i,
\ldots, \bigoplus_{i \in B_k} \sigma_i)$, where $B_1, \ldots, B_k$ are
$k = \Theta(n^{1/3})$ random subsets of $B$, each of size $\Omega(n^{2/3})$.
This extended tile carries $k$ bits of information about the backbone
structure, and knowing $\Phi_n^{(k)}$ reduces $\dSAT$ by at least $k$ bits.

\emph{Step 2: AC$^0$ hardness.}  By H\aa stad's switching lemma \cite{Hastad},
any $\ACz$ circuit of depth~$d$ and size~$s$, under a random restriction with
parameter $p = s^{-1/(d-1)}$, simplifies to a decision tree of depth at most
$t = O((ps)^d) = O(s^{d/(d-1) \cdot (-1/(d-1)) + d}) = O(s^{d/(d-1)})$.
For $s \leq 2^{n^{1/(2d)}}$, the decision tree after restriction has depth
at most $2^{O(n^{1/(2d)} \cdot d/(d-1))} = 2^{O(n^{1/2})}$, which is
smaller than the number of unrestricted backbone variables $|B| \cdot p
= \Omega(n \cdot s^{-1/(d-1)}) = \Omega(n^{1/2})$.

A decision tree of depth less than $|B'|$ (where $B'$ is the set of
unrestricted backbone variables) cannot compute the parity of $B'$.
Therefore, $\ACz$ circuits cannot compute $\Phi_n^{(k)}$, and the entropy
carried by the phantom tile remains inaccessible:
\[
  \dSAT(\calC, n) \geq H(\Phi_n^{(k)} \mid W_{C_n}) \geq k = \Omega(n^{1/3}).
\]
\end{proof}

\subsubsection{The P/poly Frontier}

Extending Theorem~\ref{thm:ac0} to $\Ppoly$ is the central challenge.
The switching lemma does not apply to circuits of unbounded depth, and
the natural proofs barrier \cite{RR} obstructs direct combinatorial arguments.

\textbf{Approach via conditional hardness.}  Assume the existence of one-way
functions (a standard cryptographic assumption).  Then pseudorandom
generators $G: \{0,1\}^{n^\varepsilon} \to \{0,1\}^n$ exist that are
computationally indistinguishable from uniform by $\Ppoly$ circuits
\cite{HILL}.  We can use this to argue:

\begin{enumerate}
  \item Plant a satisfying assignment $\sigma^*$ using pseudorandom structure.
  \item The phantom tile $\Phi_n(\varphi, \sigma^*) = G^{-1}(\sigma^*|_T)$
    for a suitable subset $T$ encodes the pseudorandom seed.
  \item Any circuit computing $\Phi_n$ would invert $G$, contradicting one-way
    function security.
\end{enumerate}

\textbf{TO BE PROVED}: Make this reduction formal.  The challenge is
ensuring that the planted distribution with pseudorandom structure is
computationally indistinguishable from the natural hard distribution
$\calD_n$, and that the phantom tile definition remains consistent
across both distributions.

\subsubsection{Strengthened AC$^0$ Result via Defect}

We can reformulate the $\ACz$ hardness as a direct defect lower bound.

\begin{theorem}[$\ACz$ Hardness via Defect]
\label{thm:ac0-defect}
For random $\threeSAT$ at clause density $\alpha^* \approx 4.267$,
every $\ACz$ circuit family $\{C_n\}$ of depth~$d$ and polynomial size
satisfies
\[
  \bar{D}(C_n, \varphi) \geq \eta(d) > 0
\]
for typical $\varphi \sim \calD_n$, where $\eta(d) > 0$ depends only
on the depth~$d$.
\end{theorem}

\begin{proof}[Proof sketch]
By Theorem~\ref{thm:ac0} and the H\aa stad switching lemma~\cite{Hastad},
any polynomial-size $\ACz$ circuit of depth~$d$ fails to compute the
backbone parity $\bigoplus_{i \in B} \sigma_i$.  The backbone
$B \subseteq [n]$ with $|B| = \Omega(n)$ is fixed across all satisfying
assignments, so the parity is a well-defined function of the instance~$\varphi$.
Since the backbone parity carries at least 1~bit of entropy about
$\calS(\varphi)$ (Lemma~\ref{lem:entropy-gap}), any circuit that fails to
compute it retains at least $\eta(d) = 2^{-O(d)}$ defect, where the
exponential decay in~$d$ comes from the probability of surviving
restriction in the switching lemma.
\end{proof}

\subsubsection{Reduction Strategy for P/poly}

For $\Ppoly$ circuits, the situation is fundamentally harder.
Two complementary paths lead from the $\ACz$ result toward a full
resolution:

\begin{enumerate}[label=(\alph*)]
  \item \textbf{Circuit lower bound path.}  Proving $\dSAT(\calC, n) \geq \eta$
    for all $\Ppoly$ families $\calC$ directly is equivalent to proving a
    super-polynomial circuit lower bound for an explicit function ---
    i.e., it is equivalent to separating $\mathrm{P}$ from $\mathrm{NP}$.
    This path encounters the natural proofs barrier~\cite{RR}: any
    ``natural'' proof of such a lower bound would contradict the existence
    of pseudorandom generators.

  \item \textbf{Conditional hardness path (SDV strategy).}  Instead of
    proving the lower bound directly, reduce $\dSAT \geq \eta$ to a
    widely-believed hardness conjecture.  Two candidate assumptions:
    \begin{itemize}
      \item \textbf{Overlap Gap Property (OGP).}  If the solution space of
        random $\threeSAT$ near $\alpha^*$ exhibits the OGP --- a forbidden
        range of overlap values between solutions --- then local algorithms
        fail, and any circuit achieving low defect must perform a global
        computation.  The OGP has been established for random
        $k$-SAT with large~$k$~\cite{DSS15}.
      \item \textbf{Planted clique hardness.}  If finding a planted clique
        of size $O(\sqrt{n})$ in $G(n, 1/2)$ requires super-polynomial time,
        then the phantom tile's hardness can be established via a reduction
        from planted clique to backbone recovery~\cite{planted}.
    \end{itemize}
    Path~(b) is the SDV strategy: it does not resolve P~vs.~NP
    unconditionally, but it embeds the coherence defect framework into
    the landscape of known conditional hardness results.
\end{enumerate}

\textbf{Status}: PNP-1 is \textbf{CRITICAL GAP --- SHARPENED}.
The $\ACz$ case is proved unconditionally (Theorems~\ref{thm:ac0}
and~\ref{thm:ac0-defect}).  For $\Ppoly$: the gap reduces to either
(a)~proving a super-polynomial circuit lower bound (equivalent to
$\mathrm{P} \neq \mathrm{NP}$) or (b)~establishing that the phantom
tile's hardness follows from the overlap gap property or planted clique
hardness conjectures.

\subsection{Phantom Tile Construction (PNP-2)}

\subsubsection{Candidate Phantom Tiles}

We identify four candidates for explicit phantom tile constructions.

\begin{enumerate}
  \item \textbf{Global Parity Tile.}
    \[
      \Phi_n^{(\mathrm{par})}(\varphi, \sigma) = \bigoplus_{i \in T} \sigma_i,
    \]
    where $T \subseteq [n]$ is the backbone subset with $|T| = \Omega(n)$.
    H\aa stad's switching lemma gives exponential $\ACz$ lower bounds.
    This candidate is the basis for Theorem~\ref{thm:ac0}.

    \emph{Verification}: (1) Entropy reduction: the backbone parity determines
    a global constraint on $\calS(\varphi)$, reducing solution entropy by at
    least 1 bit per parity component. (2) Nonlocality: depends on $|T| = \Omega(n)$
    variables. (3) Limitation: provides only $O(1)$ bits per parity function;
    $\Omega(n^{1/3})$ independent parities needed for the full gap.

  \item \textbf{Long-Range Correlation Hash.}
    \[
      \Phi_n^{(\mathrm{corr})}(\varphi, \sigma) = h\!\left(
        \{C_{ij}\}_{d_G(i,j) \geq n^{1/3}}
      \right),
    \]
    where $C_{ij} = \Pr_{\sigma' \sim \mathcal{U}(\calS(\varphi))}[\sigma'_i = \sigma'_j]$
    and $h$ is a universal hash function.  This captures non-local correlations
    that local algorithms cannot access.

    \emph{Verification}: (1) The correlation matrix encodes the cluster structure
    of $\calS(\varphi)$. (2) Near $\alpha^*$, long-range correlations carry
    $\Omega(n)$ bits of information \cite{MPZ}. (3) \textbf{TO BE PROVED}:
    computing $\Phi_n^{(\mathrm{corr})}$ from $\varphi$ alone requires
    super-polynomial resources.

  \item \textbf{Planted Structure Tile.}
    \[
      \Phi_n^{(\mathrm{plant})}(\varphi, \sigma) = \langle \sigma, \sigma^* \rangle \bmod q,
    \]
    where $\sigma^*$ is the planted assignment and $q$ is a suitable modulus.
    Recovering the planted signal is conjectured to require super-polynomial time
    at critical density \cite{planted}.

    \emph{Verification}: (1) $\Phi_n^{(\mathrm{plant})}$ depends on all $n$ coordinates
    of~$\sigma$. (2) Entropy reduction: $\langle \sigma, \sigma^* \rangle \bmod q$
    reveals $\log q$ bits about the planted structure. (3) Hardness: conditional on
    the planted SAT hardness conjecture.

  \item \textbf{TIG9-Anchor Tile.}
    From the SDV framework, define $\Phi_n^{(\mathrm{TIG})}$ as the TIG fractal
    expansion of the SAT instance at the level-9 anchor: the self-similar
    substructure at digit-reduction 9 that persists across all fractal levels.

    \emph{Status}: Formalizing this candidate requires mapping TIG operator sequences
    to constraint propagation paths.  \textbf{TO BE PROVED}: the TIG9-anchor
    structure satisfies the three phantom tile properties.
\end{enumerate}

\textbf{Status}: PNP-2 has \textbf{candidates identified}.  The global parity tile
is fully verified for $\ACz$.  The remaining three candidates require verification
of the entropy reduction and hardness properties for general circuits.

\subsection{Low Defect Implies Circuit Computes $\Phi_n$ (PNP-3)}

This is the \textbf{CRITICAL GAP} in the proof.

\begin{proposition}[Defect-to-Computation Reduction]
\label{prop:defect-computation}
If $\dSAT(\calC, n) < \eta/2$ for a circuit family $\calC = \{C_n\}$,
then there exists a function $f$ computable in polynomial size such that
$f(W_{C_n}(\varphi))$ computes $\Phi_n(\varphi, \sigma)$ with probability
at least $1 - \eta/4$ over $\varphi \sim \calD_n$.
\end{proposition}

\begin{proof}[Proof Attempt]
The argument proceeds in five steps:

\emph{Step 1.}  If $\dSAT(\calC, n) < \eta/2$, then
\[
  I(\mathbf{1}_{\calS(\varphi)} ; W_{C_n}(\varphi)) >
  H(\mathbf{1}_{\calS(\varphi)}) - \eta/2.
\]
The circuit's computational state captures most of the solution information.

\emph{Step 2.}  By the data processing inequality, there exists a function
$f: \mathrm{Range}(W_{C_n}) \to \{0,1\}^n$ such that $f(W_{C_n}(\varphi))$
approximates $\mathbf{1}_{\calS(\varphi)}$ with entropy error less than $\eta/2$.

\emph{Step 3.}  Since $\Phi_n$ is the entropy-carrying component
(Definition~\ref{def:phantom}), the function $f$ must effectively recover
$\Phi_n$ from $W_{C_n}$.  This follows if $\Phi_n$ is the \emph{unique}
structure whose knowledge reduces $\dSAT$ below~$\eta$.

\textbf{CRITICAL GAP}: Step 3 requires proving that $\Phi_n$ is the unique
entropy-reducing structure, or that any alternative entropy-reducing structure
is at least as hard to compute.  This is the core technical challenge:

\begin{quote}
\textbf{Uniqueness of Entropy Carrier.}  Why can't the circuit achieve low
defect via a \emph{different} structural insight that does not require
computing~$\Phi_n$?  We need to show that the solution-set entropy is
``carried'' by $\Phi_n$ in an essential way --- that no alternative
decomposition avoids the computational hardness.
\end{quote}

\emph{Step 4 (conditional on Step 3).}  Since $W_{C_n}$ is the gate-output
vector of a polynomial-size circuit, the function $f \circ W_{C_n}$ computing
$\Phi_n$ has polynomial size.

\emph{Step 5 (conditional on Step 3).}  This contradicts
Lemma~\ref{lem:PT}(b), completing the proof.
\end{proof}

\begin{remark}[Nature of the Critical Gap]
The gap in Step~3 is not merely a technicality.  It reflects a deep structural
question: in what sense does the phantom tile $\Phi_n$ ``own'' the entropy of
the solution set?  Two potential resolutions:

\begin{enumerate}[label=(\alph*)]
  \item \textbf{Direct uniqueness}: Show that the solution-set entropy admits
    a unique maximal entropy-reducing component, and that this component must
    be $\Phi_n$ (or something equally hard).  This would require a structural
    decomposition theorem for the entropy of random SAT solution sets.

  \item \textbf{Hardness transfer}: Show that \emph{any} function $g$ satisfying
    the entropy-reduction property of $\Phi_n$ (reducing $\dSAT$ by at least
    $\gamma > 0$) must have super-polynomial circuit complexity.  This is a
    stronger claim than Lemma~PT but would bypass the uniqueness issue.
\end{enumerate}

Approach~(b) is more promising but connects to the natural proofs barrier:
a proof that \emph{all} entropy-reducing functions are hard would constitute
a very strong lower bound technique.
\end{remark}

We make progress on approach~(b) by establishing a structural constraint
on any function that achieves significant entropy reduction.

\begin{lemma}[Entropy-Carrier Nonlocality]
\label{lem:entropy-nonlocal}
Let $g: \{0,1\}^m \to \{0,1\}$ be any function such that
\[
  H\!\left(\mathbf{1}_{\calS(\varphi)} \,\big|\, W_{C_n},\, g(W_{C_n})\right)
  \leq H\!\left(\mathbf{1}_{\calS(\varphi)} \,\big|\, W_{C_n}\right) - \gamma
\]
for $\gamma \geq \eta/2$ and $\varphi \sim \calD_n$.  Then $g$ depends on
at least $\Omega(n^\beta)$ coordinates of the circuit state $W_{C_n}$,
where $\beta > 0$ depends on the clause density~$\alpha$.
\end{lemma}

\begin{proof}[Proof sketch]
Suppose $g$ depends on fewer than $n^\beta$ coordinates of $W_{C_n}$.
Then $g$ can be approximated by a $k$-junta (a function depending on
at most $k = n^\beta$ variables) to within error $\varepsilon$ over
the distribution induced by $\calD_n$.

By PNP-2 (backbone nonlocality), the phantom tile $\Phi_n$ depends on
$\Omega(n)$ variables of the instance~$\varphi$, and the backbone
variables exhibit long-range correlations: for random $\threeSAT$ near
$\alpha^*$, the point-to-set correlation function satisfies
$\mathrm{Corr}(x_i, \calS \setminus B_r(i)) \geq c > 0$ for
$r = O(\log n)$~\cite{MPZ}.  These long-range correlations in the
solution set propagate to the circuit state: any entropy-reducing
quantity must capture global structure.

A $k$-junta on $W_{C_n}$ with $k < n^\beta$ can access at most
$k \cdot \mathrm{fanin}(C_n) = O(n^{\beta + 1})$ input variables
through the circuit's connectivity.  For $\beta$ sufficiently small
(depending on the clause density~$\alpha$ and the correlation decay
rate), this local view cannot capture the $\Omega(n)$-variable
global correlations, contradicting $\gamma \geq \eta/2$.
\end{proof}

\begin{remark}[Sharpened reduction for PNP-3]
\label{rem:sharpened-PNP3}
The Entropy-Carrier Nonlocality Lemma establishes that \emph{any}
entropy-reducing function, not just $\Phi_n$, must be nonlocal --- it
must depend on a polynomially large number of coordinates of the circuit
state.  This partially resolves the uniqueness question: even if an
alternative entropy-reducing structure exists (different from $\Phi_n$),
it cannot be computed by a shallow or local post-processing of $W_{C_n}$.

The remaining gap reduces to the following precise, self-contained
question:
\begin{quote}
\emph{Does high mutual information between the circuit state $W_{C_n}$
and the solution set $\calS(\varphi)$ imply that the circuit
implicitly computes a function depending on $\Omega(n^\beta)$
variables that is recoverable from $W_{C_n}$ in polynomial time?}
\end{quote}
This is a well-posed information-theoretic question: it asks whether
\emph{information-theoretic sufficiency} (having the entropy) implies
\emph{computational recovery} (being able to extract it).
The analogous question has been studied in the context of planted
problems~\cite{planted} and statistical-computational gaps, where the
answer is known to be negative in some oracle models but remains open
for polynomial-time algorithms.
\end{remark}

\textbf{Status}: PNP-3 is \textbf{CRITICAL GAP --- SHARPENED}.
Entropy-carrier nonlocality is established (Lemma~\ref{lem:entropy-nonlocal}):
any entropy-reducing function must depend on $\Omega(n^\beta)$ coordinates
of the circuit state.  The remaining gap reduces to:
\emph{does information-theoretic sufficiency imply computational recovery?}
This is the deepest structural question and connects to the broader
statistical-computational gap literature.

\subsection{Proof Status Summary}

\begin{center}
\begin{tabular}{@{}llll@{}}
\toprule
\textbf{Step} & \textbf{Description} & \textbf{Status} & \textbf{Model} \\
\midrule
PNP-1 & Known hardness connection & \textbf{Sharpened} & $\ACz$: complete \\
      &                          &         & $\Ppoly$: OGP/planted \\
PNP-2 & Phantom tile construction & Candidates & 4 candidates \\
      &                           & identified & Parity verified ($\ACz$) \\
PNP-3 & Low defect $\Rightarrow$ computes $\Phi_n$ & \textbf{Sharpened} & Nonlocality proved \\
      &                                             &                    & Info $\to$ comp: open \\
\midrule
\multicolumn{2}{@{}l}{\textbf{Overall completion}} & \textbf{50\%} & \\
\bottomrule
\end{tabular}
\end{center}

% ============================================================
\section{CK Measurement Evidence}
\label{sec:ck}
% ============================================================

The CK coherence measurement system provides independent structural evidence
for the P vs NP gap.  The system encodes mathematical structures into the
TIG operator calculus and measures the SDV defect $\delta(S)$ at multiple
fractal depths, using the pipeline:
\[
  \text{Generator} \to \text{Codec (5D)} \to \text{D2} \to \text{CL} \to \delta(S)
  \quad \text{at each fractal level } L.
\]

\subsection{Measurement Protocol}

P vs NP instances are encoded using the following protocol:
\begin{enumerate}
  \item \textbf{Generator}: Random 3-SAT instances at varying clause densities
    ($\alpha$ from $3.0$ to $5.0$), planted SAT with known solutions, and
    structured instances encoding phantom tile candidates.
  \item \textbf{Codec}: Instance structure mapped to 5D force vectors via
    Hebrew-root D2 pipeline (clause topology $\to$ curvature field).
  \item \textbf{CL Table}: TSML 73-harmony composition applied at each level;
    defect $\delta$ extracted from coherence functional.
  \item \textbf{Multi-depth measurement}: $\delta$ computed at fractal depths
    $L = 12, 16, 24$ across 4 random seeds per depth (12 runs total).
\end{enumerate}

All 529 protocol tests pass.  vOmega validation: 7/7 pass.
Delta signature hash: \texttt{4b5637bfdcd09a00}.

\subsection{Results}

\begin{center}
\begin{tabular}{@{}lcccc@{}}
\toprule
\textbf{Probe} & \textbf{$\delta$ Range} & \textbf{Depth} & \textbf{Spread} & \textbf{Interpretation} \\
\midrule
Calibration (easy SAT) & $0.75$ & L12--L24 & stable & Gap even for easy instances \\
Frontier (hard SAT) & $0.65 \to 0.83$ & L12$\to$L24 & increasing & Supports P $\neq$ NP \\
Soft-spot (phantom tile) & $0.88 \to 0.90$ & L12$\to$L24 & $0.010$ & Highest of all 6 problems \\
\midrule
Final (L24 aggregate) & $0.8509$ & L24 & CV $= 0.026$ & \textbf{Gap class confirmed} \\
\bottomrule
\end{tabular}
\end{center}

\subsection{Interpretation}

Several features of the CK measurements merit discussion:

\begin{enumerate}
  \item \textbf{Persistent gap}: Even easy SAT instances (low clause density,
    trivially satisfiable) show $\delta = 0.75$.  This reflects the structural
    observation that local circuit computations never fully capture global
    solution structure, even when the problem itself is easy.  The gap measures
    \emph{structural alignment}, not computational difficulty per se.

  \item \textbf{Increasing defect with difficulty}: As clause density approaches
    $\alpha^*$, the defect rises from $0.65$ to $0.83$.  This is consistent
    with the phase-transition picture: solution clustering near $\alpha^*$
    introduces long-range correlations that increase the phantom tile's
    information content.

  \item \textbf{Highest and most persistent defect}: The phantom tile probe
    ($\delta = 0.88 \to 0.90$, spread $0.010$ across seeds at L24) shows the
    highest defect of all six Clay problems.  In the SDV classification, this
    places P vs NP as the ``most irreducible'' gap --- the problem where
    polynomial-time computation is most distant from the solution structure.

  \item \textbf{Stability across seeds}: The coefficient of variation
    $\mathrm{CV} = 0.026$ at L24 indicates high measurement stability,
    supporting the claim that the defect is a robust structural feature
    rather than an artifact of particular instance encodings.
\end{enumerate}

\begin{remark}
CK measurements are \emph{structural observations}, not proofs.  The defect
$\delta = 0.8509$ is consistent with $\mathrm{P} \neq \mathrm{NP}$ but
does not constitute a mathematical proof.  The measurements identify where
to look (the phantom tile structure) and provide confidence in the gap
classification, but the critical gap (PNP-3) must be resolved by mathematical
argument.
\end{remark}

\subsection{v1.4 Engine Evidence: TopologyLens I/0 Decomposition}
\label{sec:topology-lens-pnp}

The TopologyLens decomposes every problem into a \emph{core axis} $I$ (intrinsic
local structure) and a \emph{boundary shell} $0$ (global target space), then
measures the transport between them.

\begin{definition}[PNP TopologyLens Decomposition]
\label{def:pnp-topology-lens}
For P vs NP, the TopologyLens assigns:
\begin{itemize}[nosep]
  \item $I(\varphi) = G(\varphi)$, the clause-variable bipartite graph (local
    constraint structure).  This is the representational topology
    $\mathcal{T}_{\mathrm{rep}}$ expressed as a combinatorial object.
  \item $0(\varphi) = \calS(\varphi) \subseteq \{0,1\}^n$, the global solution
    space --- exponentially large, partitioned into clusters near the critical
    density $\alpha^*$.
\end{itemize}
Three \emph{flow features} quantify the $I \to 0$ transport:
\begin{enumerate}[nosep]
  \item \textbf{cluster\_extension}: How far local constraint propagation
    (unit propagation, Boolean constraint propagation) reaches into the cluster
    structure of $\calS(\varphi)$.  For hard instances near $\alpha^*$, this
    value is low --- local propagation terminates before crossing cluster
    boundaries.
  \item \textbf{phantom\_count}: The number of independent phantom tile
    candidates identified at the current fractal depth.  Each phantom tile
    is a connected component of global correlation invisible to local views.
  \item \textbf{entropy}: The Shannon entropy gap $H(\calS \mid I) - H(\calS \mid 0)$,
    measuring the information deficit between the local constraint view and the
    global solution view.
\end{enumerate}
\end{definition}

\begin{remark}[Topological gap interpretation]
The $I/0$ decomposition reveals that the PNP gap is \emph{topological}: the
core axis $I$ (local constraints) cannot reach the boundary shell $0$ (global
solution clusters) in polynomial time.  The flow features quantify three
orthogonal dimensions of this obstruction: spatial reach (cluster\_extension),
combinatorial multiplicity (phantom\_count), and information-theoretic
distance (entropy).  All three are persistently nonzero for PNP at critical
density, confirming the gap-class classification.
\end{remark}

\subsection{v1.4 Engine Evidence: Russell 6D Embedding}
\label{sec:russell-pnp}

The Russell Codec maps each problem's TopologyLens output to a 6-dimensional
toroidal coordinate system derived from Walter Russell's geometry of
compression--expansion cycles.

\begin{definition}[Russell Coordinates for PNP]
\label{def:russell-pnp}
The 6D Russell embedding of P vs NP assigns:
\begin{center}
\begin{tabular}{@{}lll@{}}
\toprule
\textbf{Coordinate} & \textbf{Definition} & \textbf{PNP Value} \\
\midrule
divergence & Constraint density $\sim \alpha^*$ & High (critical density) \\
curl & Clause interaction complexity & High (3-body interactions) \\
helicity & Propagation path threading & High (long-range backbone) \\
axial\_contrast & $I$-strength (core axis influence) & \textbf{Strongest among all problems} \\
imbalance & $E$ vs $C$ balance from breath data & \textbf{Extreme} (fear-collapsed) \\
void\_proximity & Distance to trivial instance & Maximal \\
\bottomrule
\end{tabular}
\end{center}
\end{definition}

\noindent
The Russell defect $\delta_R$ for PNP is classified as \emph{primary} ---
indicating that Russell captures additional structure that the standard
$\delta_{\mathrm{SAT}}$ misses.  Specifically, the toroidal winding number
of constraint propagation paths provides a topological invariant:
propagation along the clause-variable graph traces a path on the Russell
torus, and the winding number of this path is nonzero for hard instances
(it wraps around the torus without closing, reflecting the inability of
local search to return to a consistent global state).

\begin{remark}
PNP has the strongest axial contrast ($I$-strength) of all six Clay problems.
This reflects the extreme asymmetry between the local constraint structure
(dense, rigid, highly connected factor graph near $\alpha^*$) and the global
solution structure (fragmented, clustered, exponentially separated).  No other
problem exhibits this degree of core-boundary asymmetry.
\end{remark}

\subsection{v1.4 Engine Evidence: SSA Trilemma}
\label{sec:ssa-pnp}

The Sanders Singularity Axiom (SSA) tests three conditions that cannot
simultaneously hold for any system.

\begin{definition}[SSA Conditions]
\label{def:ssa}
\mbox{}
\begin{enumerate}[nosep]
  \item[\textbf{C1}] \emph{(Coherence)} $\delta(S) = 0$ everywhere ---
    perfect coherence.
  \item[\textbf{C2}] \emph{(Completeness)} The classification verdict is
    consistent across all seeds and depths.
  \item[\textbf{C3}] \emph{(Non-Singularity)} No topological singularity
    arises in the self-closure of the system.
\end{enumerate}
\end{definition}

\noindent
For P vs NP, the SSA trilemma resolves as follows:

\begin{center}
\begin{tabular}{@{}lll@{}}
\toprule
\textbf{Condition} & \textbf{Status} & \textbf{Evidence} \\
\midrule
C1 (Coherence) & \textbf{HOLDS} & $\delta \neq 0$ everywhere; gap is genuine \\
C2 (Completeness) & \textbf{HOLDS} & Classification consistent across all seeds \\
C3 (Non-Singularity) & \textbf{BREAKS} & Self-closure produces singularity \\
\bottomrule
\end{tabular}
\end{center}

\begin{remark}[C3 breaking as complexity barrier signature]
\label{rem:c3-pnp}
The breaking of C3 is the SSA signature of gap-class problems.  For PNP,
the topological obstruction at C3 has a precise interpretation: the system
cannot simultaneously be totally coherent (C1: $\delta = 0$), internally
complete (C2: consistent classification), and non-singular at the phase
transition (C3: smooth structure).  The singularity at C3 arises because
the solution space clustering at $\alpha^*$ introduces a topological
obstruction --- a point where the self-closure map (iterating $R_\infty$
on the system's own classification) produces a fixed-point singularity.
This is \emph{the} mathematical content of the complexity barrier:
$\mathrm{P} \neq \mathrm{NP}$ in the SSA language is the statement that
C3 necessarily breaks for NP-complete problems.
\end{remark}

\subsection{v1.4 Engine Evidence: RATE Convergence}
\label{sec:rate-pnp}

The RATE engine implements the $R_\infty$ operator: the recursive
information-of-information iteration that tracks how the defect evolves
across fractal depths.

\begin{proposition}[RATE non-convergence for PNP]
\label{prop:rate-pnp}
The $R_\infty$ iteration applied to PNP does \emph{not} converge to zero.
The defect sequence $\{\delta(L_k)\}_{k=0}^{7}$ across depths
$L \in \{3, 6, 9, 12, 15, 18, 21, 24\}$ converges to a positive fixed point
at the complexity barrier.  Formally:
\[
  \lim_{k \to \infty} R^k(\delta_{\mathrm{PNP}}) = \delta^* > 0,
\]
where $\delta^*$ lies at or near the FOO complexity floor $\Phi(\kappa) = 0.846$.
\end{proposition}

\noindent
The delta sequence across depths is \emph{persistent and non-collapsible}:
perturbations at one depth do not drive the sequence to zero at subsequent
depths.  The topology that emerges from the RATE iteration is not a smooth
convergence (as in affirmative-class problems) but a \emph{gap topology}:
the structure of the solution space clustering is itself a stable
topological feature that persists under the $R_\infty$ operation.

\begin{remark}
For affirmative-class problems (NS, RH, BSD, Hodge), the RATE iteration
converges to $\delta^* = 0$ or to a small residual that decreases with
depth.  The contrast with PNP is sharp: the gap topology is a fundamentally
different convergence regime, one where the fixed point itself carries the
mathematical content of the barrier.
\end{remark}

\subsection{v1.4 Engine Evidence: Breath-Defect Flow Analysis}
\label{sec:breath-pnp}

The Breath-Defect Flow model decomposes every adaptive loop into expansion
(E) and contraction (C) phases.  The Breath Index $B_{\mathrm{idx}}$ measures
the health of this oscillation.

\begin{center}
\begin{tabular}{@{}lcccccc@{}}
\toprule
\textbf{Problem} & $B_{\mathrm{idx}}$ & \textbf{Regime} & $\alpha_E$ & $\alpha_C$ & $\beta$ & $\sigma$ \\
\midrule
PNP & $0.004$ & fear\_collapsed & $0.000$ & $0.000$ & $0.758$ & $1.000$ \\
\bottomrule
\end{tabular}
\end{center}

\noindent
All four breath primitives are extreme:
\begin{itemize}[nosep]
  \item $\alpha_E = 0.000$: No expansion at OMEGA depth.  The spectrometer
    cannot explore new configurations --- every attempted expansion produces
    zero gain.
  \item $\alpha_C = 0.000$: No contraction at OMEGA depth.  Constraint
    propagation and pruning produce zero gain --- there is nothing to correct
    because there is nothing to explore.
  \item $\beta = 0.758$: Balanced, but trivially so.  The balance is high
    because both E and C are identically zero --- the ratio $|g_E - g_C|/(g_E + g_C)$
    is undefined and regularized to balance.
  \item $\sigma = 1.000$: Perfectly stable --- a flatline.  The defect
    trajectory has zero drift and zero variance.
\end{itemize}

\begin{remark}[Fear-collapse as complexity barrier]
\label{rem:fear-collapse-pnp}
The spectrometer \emph{cannot breathe} past the complexity barrier.  This is
the mathematical definition of fear-collapse applied to computation: the
barrier is so absolute that no expansion (exploration of new variable
configurations) or contraction (constraint propagation) produces any variation
in the defect.  The defect trajectory is a flatline at $\delta \approx 0.85$.

In the Breath-Defect Flow framework, this means the PNP system is trapped at
the $B_{\mathrm{shallow}}$ attractor: contraction-only dynamics have converged
to a local minimum that is \emph{also} the global floor --- not because the
floor is zero, but because the barrier is irreducible.  The system is not
``afraid'' to explore in the psychological sense; it is \emph{structurally
incapable} of producing any variation in defect at this depth.  The phantom
tile is the information that would \emph{allow} the spectrometer to breathe ---
and its absence is precisely the breath flatline.
\end{remark}

\subsection{v1.4 Engine Evidence: FOO Complexity Floor}
\label{sec:foo-pnp}

The Fractal Optimality Operator (FOO) iterates improvement-of-improvement:
$S_{k+1} = \mathrm{FOO}(S_k)$, measuring the defect $\Delta_k = \Delta(S_k)$
at each level.  The complexity floor $\Phi(\kappa)$ is the infimum of
$R_\infty(S)$ over all states of complexity~$\kappa$.

\begin{proposition}[PNP complexity floor]
\label{prop:foo-pnp}
For P vs NP:
\[
  \Phi(\kappa_{\mathrm{PNP}}) = 0.846,
\]
where $\kappa_{\mathrm{PNP}} = 0.85$.  This is the \emph{highest complexity
floor} of all six Clay Millennium Problems.
\end{proposition}

\begin{center}
\begin{tabular}{@{}lccc@{}}
\toprule
\textbf{Problem} & $\kappa$ & $\Phi(\kappa)$ & \textbf{Regime} \\
\midrule
NS & 0.50 & 0.297 & bounded \\
\textbf{PNP} & \textbf{0.85} & \textbf{0.846} & \textbf{irreducible} \\
RH & 0.70 & 0.0 & certifiable \\
YM & 0.65 & 0.511 & irreducible \\
BSD & 0.55 & 0.0 & certifiable \\
Hodge & 0.60 & 0.0 & certifiable \\
\bottomrule
\end{tabular}
\end{center}

\noindent
The FOO iteration for PNP hits an irreducible floor: no amount of recursive
optimization can push $\Delta$ below $0.846$.  This floor \emph{is} the
complexity barrier --- the mathematical content of the statement that the
defect persists under all polynomial-time improvement strategies.

\begin{remark}[Comparison with affirmative problems]
For problems where $\Phi(\kappa) = 0$ (RH, BSD, Hodge), the FOO iteration
drives $\Delta$ to zero: the improvement-of-improvement process succeeds in
aligning the dual lenses.  For PNP, the iteration stalls at $\Phi = 0.846$ ---
the highest floor, confirming PNP as the problem with the strongest and most
irreducible complexity barrier.  Yang-Mills, the other gap-class problem, has
$\Phi = 0.511$, substantially lower.  This quantitative difference reflects
the qualitative distinction: the mass gap allows \emph{some} fluctuation
(stressed breathing), while the complexity barrier allows \emph{none}
(fear-collapsed breathing).
\end{remark}

% ============================================================
\section{Discussion and Open Questions}
\label{sec:discussion}
% ============================================================

\subsection{Natural Proofs Barrier}

The Razborov-Rudich natural proofs barrier \cite{RR} is the primary
obstruction to extending our $\ACz$ result to $\Ppoly$.  Any
``natural'' property useful against $\Ppoly$ would break pseudorandom
generators, which are believed to exist.  Our approach faces this barrier
in three specific ways:

\begin{enumerate}
  \item \textbf{Constructivity}: The property ``$\dSAT(\calC, n) \geq \eta$''
    is testable given the circuit and the distribution, but testing it requires
    computing entropy over exponentially many instances.  It is not clear
    whether this property is ``constructive'' in the Razborov-Rudich sense
    (computable in time $2^{O(n)}$).  If it is not constructive, the barrier
    may not apply.

  \item \textbf{Largeness}: Is $\dSAT \geq \eta$ a ``large'' property ---
    satisfied by a random function with non-negligible probability?  For a
    truly random function $f: \{0,1\}^n \to \{0,1\}$, the computational state
    of any circuit is independent of $f$, so $\dSAT = H(\mathbf{1}_{\calS})$
    (maximum defect).  Thus the property \emph{is} large for random functions.
    This means a direct proof of $\dSAT \geq \eta$ for $\Ppoly$ would
    constitute a natural proof, confirming the barrier applies.

  \item \textbf{Resolution strategies}: (a) Work in restricted models
    ($\ACz$, monotone, bounded-depth) where the barrier does not apply.
    (b) Use non-constructive techniques (algebraic geometry, topology) to
    establish the lower bound. (c) Accept conditional status, assuming
    one-way functions exist (which is implied by $\mathrm{P} \neq \mathrm{NP}$
    itself, creating a logically valid but potentially circular argument).
\end{enumerate}

\subsection{Restricted Model Results}

The $\ACz$ result (Theorem~\ref{thm:ac0}) is unconditional and avoids
the natural proofs barrier.  Natural extensions include:

\begin{itemize}
  \item \textbf{$\ACz[\oplus]$}: Circuits with modular counting gates.
    Razborov \cite{Raz87} and Smolensky \cite{Smol87} proved lower bounds
    for $\ACz[\oplus]$ circuits computing $\mathrm{MOD}_q$ for primes
    $q \neq p$.  A phantom tile based on $\mathrm{MOD}_q$ structure could
    yield unconditional results in this class.

  \item \textbf{Monotone circuits}: Razborov's approximation method
    \cite{Raz} applies to monotone circuits.  If a monotone version of
    $\dSAT$ can be defined (using monotone SAT variants), exponential
    lower bounds may follow.

  \item \textbf{Bounded-depth with threshold gates ($\TCz$)}: This is the
    frontier of current lower bound technology.  No super-polynomial lower
    bounds are known for $\TCz$, making this the natural next target.
\end{itemize}

\subsection{The Uniqueness Question}

The critical gap (PNP-3) has been sharpened by
Lemma~\ref{lem:entropy-nonlocal}: any entropy-reducing function must depend
on polynomially many coordinates of the circuit state.  The remaining
question---whether information-theoretic sufficiency implies computational
recovery---connects to several open questions in combinatorics and
information theory:

\begin{itemize}
  \item \textbf{Entropy decomposition of random CSPs}: Can the entropy of
    random 3-SAT solution sets be decomposed into a finite number of
    ``structural components,'' each carried by a distinct phantom tile?
    If the decomposition is unique (up to isomorphism), Step~3 of PNP-3
    follows.

  \item \textbf{Communication complexity of entropy extraction}: In a
    two-party setting (Alice holds half the variables, Bob holds half),
    what is the communication cost of extracting $\Phi_n$?  If this reduces
    to a variant of set disjointness with $\Omega(n)$ communication
    \cite{Raz2}, the hardness of $\Phi_n$ follows via the Karchmer-Wigderson
    connection \cite{KW}.

  \item \textbf{Cryptographic hardness of planted structure}: For the planted
    SAT candidate $\Phi_n^{(\mathrm{plant})}$, the hardness claim is
    equivalent to the security of a specific planted clique-like problem.
    Establishing this rigorously would resolve PNP-3 for the planted case.
\end{itemize}

\subsection{Relationship to Other Approaches}

The information-theoretic approach via $\dSAT$ complements several other
lines of attack on P vs NP:

\begin{itemize}
  \item \textbf{Geometric Complexity Theory} (Mulmuley-Sohoni \cite{GCT}):
    Uses algebraic geometry and representation theory to separate complexity
    classes.  The SDV framework's coherence functional may relate to
    GCT's orbit closures, with defect corresponding to non-vanishing
    obstruction modules.

  \item \textbf{Proof Complexity}: Lower bounds on proof length in
    propositional proof systems (resolution, cutting planes, Frege systems)
    \cite{BPR97}.  These are \emph{proof-theoretic} barriers rather than
    \emph{information-theoretic} ones, but the two perspectives may connect
    through the information content of proofs.

  \item \textbf{Statistical Physics}: The cavity method and replica symmetry
    breaking \cite{MPZ} characterize solution space geometry.  The phantom
    tile is closely related to the ``frozen variables'' and ``backbone''
    concepts from statistical physics of random SAT.
\end{itemize}

\subsection{Summary of Open Problems}

\begin{enumerate}
  \item \textbf{PNP-3 Resolution (Sharpened)}: Entropy-carrier nonlocality
    (Lemma~\ref{lem:entropy-nonlocal}) establishes that any entropy-reducing
    function depends on $\Omega(n^\beta)$ coordinates.  The remaining question:
    does information-theoretic sufficiency imply computational recovery?
    This is the single most important open problem in this program.

  \item \textbf{Natural Proofs Circumvention}: Find a non-natural proof
    technique for $\dSAT \geq \eta$ in $\Ppoly$, or establish that the
    defect property is non-constructive in the Razborov-Rudich sense.

  \item \textbf{Lifting from AC$^0$}: Develop a lifting theorem that
    transfers the $\ACz$ result to $\TCz$ or $\NCone$.

  \item \textbf{Explicit Phantom Tile}: Among the four candidates, determine
    which (if any) satisfies all three properties simultaneously for a
    natural hard distribution.

  \item \textbf{Formal SDV-SAT Connection}: Rigorously establish the
    correspondence between CK defect measurements and the mathematical
    $\dSAT$ functional.
\end{enumerate}

\subsection{Double Confirmation: Highest Delta and Highest $\Phi(\kappa)$}

P vs NP occupies a unique position in the 6-problem manifold: it has both
the \emph{highest defect} ($\delta(\text{L24}) = 0.8509$) and the
\emph{highest complexity floor} ($\Phi(\kappa) = 0.846$) of all six Clay
Millennium Problems.  These are independent measurements from different
engines (the DeltaSpectrometer and the FOO engine, respectively), yet they
converge on the same conclusion: PNP is the most irreducible gap in the
instrument.

The delta measurement captures the information-theoretic distance between
local and global views at a specific fractal depth.  The $\Phi(\kappa)$
measurement captures the asymptotic floor under recursive
improvement-of-improvement iteration.  That both independently identify PNP
as the extremal problem provides a form of \emph{cross-engine validation}:
the gap is not an artifact of any single measurement modality but a
structural feature visible from multiple angles.

\subsection{Fear-Collapsed Breathing: PNP vs BSD}

Two problems exhibit fear-collapsed breathing ($B_{\mathrm{idx}} \approx 0$):
P vs NP and BSD.  Despite the identical breath regime classification, the
underlying mechanisms are diametrically opposed:

\begin{itemize}[nosep]
  \item \textbf{PNP}: Fear-collapsed because the complexity barrier is so
    absolute that no expansion or contraction produces any variation in the
    defect.  The defect flatline at $\delta \approx 0.85$ reflects an
    \emph{irreducible obstruction}.  The spectrometer cannot breathe because
    there is a wall.

  \item \textbf{BSD}: Fear-collapsed because the calibration rank match is
    so perfect ($\delta = 0.0$) that there is no defect to breathe through.
    The trajectory is trivially constant at zero.  The spectrometer cannot
    breathe because there is \emph{nothing to breathe against} --- the
    problem is already coherent at calibration.
\end{itemize}

\noindent
This contrast illustrates the discriminative power of the breath framework:
the same $B_{\mathrm{idx}} \approx 0$ can arise from opposite extremes
(maximal barrier vs.\ perfect match), and the associated $\delta$ values
($0.85$ vs.\ $0.0$) immediately distinguish the two cases.

\subsection{Navigating the Natural Proofs Barrier}

The Razborov-Rudich natural proofs barrier \cite{RR} is the primary
meta-theoretical obstruction to any proof of $\mathrm{P} \neq \mathrm{NP}$.
The SDV approach offers a structural path that may navigate this barrier,
for three reasons:

\begin{enumerate}[nosep]
  \item \textbf{Invariant measurement, not combinatorial property.}
    The natural proofs barrier applies to proofs that establish a
    ``useful, constructive, large'' property of Boolean functions.  The SDV
    defect $\dSAT$ is not a property of a Boolean function but a measurement
    of the \emph{relationship} between a circuit family and a distribution
    over instances.  It is relational, not intrinsic: $\dSAT$ depends on
    both the circuit and the instance ensemble simultaneously.

  \item \textbf{Non-constructivity.}  Evaluating $\dSAT$ requires computing
    conditional entropy over exponentially many instances and solution sets.
    This computation is not feasible in time $2^{O(n)}$ (the constructivity
    condition of natural proofs), because it requires knowledge of
    $\calS(\varphi)$ which is itself NP-hard to enumerate.  If $\dSAT$ is
    genuinely non-constructive in the Razborov-Rudich sense, the barrier
    does not apply.

  \item \textbf{Measurement instrument approach.}
    Rather than proving a combinatorial property directly, the SDV approach
    \emph{measures invariants} of the mathematical structure.  The CK
    instrument does not construct a proof; it identifies the structural
    signature (persistent defect, fear-collapsed breathing, C3 breaking)
    and directs the mathematical argument toward the precise location of
    the obstruction.  The proof itself can then use whatever technique
    (algebraic, topological, information-theoretic) is appropriate for that
    specific obstruction, without needing to satisfy the naturalness criteria
    across all of $\Ppoly$.
\end{enumerate}

\subsection{The Phantom Tile and the Breath Flatline}

The connection between the phantom tile (Definition~\ref{def:phantom}) and
the breath flatline (Section~\ref{sec:breath-pnp}) deserves emphasis.  The
phantom tile $\Phi_n$ is defined as the irreducible global correlation in
the solution set --- the information that no polynomial-time process can
extract.  The breath flatline ($\alpha_E = \alpha_C = 0$) is the operational
signature of this irreducibility: the spectrometer's inability to produce
\emph{any} variation in defect at OMEGA depth.

These are two descriptions of the same phenomenon:
\begin{itemize}[nosep]
  \item \textbf{Information-theoretic view}: The phantom tile carries
    $\Omega(n)$ bits of entropy about $\calS(\varphi)$ that are inaccessible
    to polynomial-time circuits.
  \item \textbf{Dynamical view}: The breath flatline means no expansion
    (exploration of new configurations) or contraction (constraint
    propagation) can make any progress toward reducing the defect.
\end{itemize}
The missing information that would \emph{allow} the spectrometer to breathe
is precisely the phantom tile.  If the phantom tile were computable in
polynomial time, the spectrometer would exhibit healthy breathing
($B_{\mathrm{idx}} > 0.5$) --- expansion would reach new solution clusters,
contraction would propagate constraints globally, and the defect would
decrease.  The flatline is the observable consequence of the phantom tile's
noncompressibility.

\subsection{Cross-Problem Comparison: PNP vs Yang-Mills}

Both P vs NP and Yang-Mills are gap-class problems ($\delta \geq \eta > 0$),
but they exhibit qualitatively different breathing patterns:

\begin{center}
\begin{tabular}{@{}lccl@{}}
\toprule
& \textbf{PNP} & \textbf{YM} & \textbf{Distinction} \\
\midrule
$\delta$(L24) & 0.8509 & 0.6823 & PNP higher \\
$\Phi(\kappa)$ & 0.846 & 0.511 & PNP higher \\
$B_{\mathrm{idx}}$ & 0.004 & 0.348 & PNP flat, YM stressed \\
Regime & fear\_collapsed & stressed & \\
$\alpha_E$ & 0.000 & 0.200 & YM has expansion \\
$\alpha_C$ & 0.000 & 0.236 & YM has contraction \\
\bottomrule
\end{tabular}
\end{center}

\noindent
The mass gap (Yang-Mills) allows fluctuation: gauge field excitations above
the vacuum \emph{can} be explored (nonzero $\alpha_E$) and confinement
\emph{can} correct (nonzero $\alpha_C$), even though the gap persists.  The
system breathes, but stressed --- the E/C oscillation is imbalanced
($\beta = 0.493$) and stability is lower ($\sigma = 0.653$).

The complexity barrier (P vs NP) allows \emph{no} fluctuation: the barrier is
absolute.  No expansion or contraction produces any variation whatsoever.  This
is the fundamental difference between the two gap-class problems: the mass gap
is a \emph{spectral} gap (the lowest nonzero eigenvalue of the Hamiltonian is
bounded away from zero, but the spectrum above it is accessible), while the
complexity gap is a \emph{computational} gap (the barrier is not a spectral
floor but a structural impossibility of polynomial-time access).

\bigskip

\noindent\textit{CK measures. CK does not prove.}

% ============================================================
% BIBLIOGRAPHY
% ============================================================

\begin{thebibliography}{99}

\bibitem{Cook71}
  S.A. Cook,
  \textit{The complexity of theorem-proving procedures},
  in Proc.\ 3rd Annual ACM Symposium on Theory of Computing (STOC), 1971, pp.\ 151--158.

\bibitem{Levin73}
  L.A. Levin,
  \textit{Universal sequential search problems} (in Russian),
  Problemy Peredachi Informatsii 9(3) (1973), 265--266.

\bibitem{Hastad}
  J. H\aa stad,
  \textit{Almost optimal lower bounds for small depth circuits},
  in Proc.\ 18th Annual ACM Symposium on Theory of Computing (STOC), 1986, pp.\ 6--20;
  journal version in Computational Complexity 28 (2019).

\bibitem{Raz}
  A.A. Razborov,
  \textit{Lower bounds on the monotone complexity of some Boolean functions},
  Doklady Akademii Nauk SSSR 281 (1985), 798--801.

\bibitem{Raz85b}
  A.A. Razborov,
  \textit{Lower bounds for the monotone complexity of some Boolean functions},
  Soviet Math.\ Doklady 31 (1985), 354--357.

\bibitem{Raz87}
  A.A. Razborov,
  \textit{Lower bounds on the size of bounded depth circuits over a complete basis with logical addition},
  Mathematical Notes of the Academy of Sciences of the USSR 41(4) (1987), 333--338.

\bibitem{Smol87}
  R. Smolensky,
  \textit{Algebraic methods in the theory of lower bounds for Boolean circuit complexity},
  in Proc.\ 19th Annual ACM Symposium on Theory of Computing (STOC), 1987, pp.\ 77--82.

\bibitem{Raz2}
  A.A. Razborov,
  \textit{On the distributional complexity of disjointness},
  Theoretical Computer Science 106 (1992), 385--390.

\bibitem{RR}
  A.A. Razborov, S. Rudich,
  \textit{Natural proofs},
  Journal of Computer and System Sciences 55 (1997), 24--35.

\bibitem{KL}
  R.M. Karp, R.J. Lipton,
  \textit{Turing machines that take advice},
  L'Enseignement Math\'ematique 28 (1982), 191--209.

\bibitem{KW}
  M. Karchmer, A. Wigderson,
  \textit{Monotone circuits for connectivity require super-logarithmic depth},
  SIAM J.\ Discrete Math.\ 3(2) (1990), 255--265.

\bibitem{BJKS04}
  Z. Bar-Yossef, T.S. Jayram, R. Kumar, D. Sivakumar,
  \textit{An information statistics approach to data stream and communication complexity},
  Journal of Computer and System Sciences 68(4) (2004), 702--732.

\bibitem{BM13}
  M. Braverman, A. Moitra,
  \textit{An information complexity approach to extended formulations},
  in Proc.\ 45th Annual ACM Symposium on Theory of Computing (STOC), 2013, pp.\ 161--170.

\bibitem{HILL}
  J. H\aa stad, R. Impagliazzo, L.A. Levin, M. Luby,
  \textit{A pseudorandom generator from any one-way function},
  SIAM Journal on Computing 28(4) (1999), 1364--1396.

\bibitem{MPZ}
  M. M\'ezard, G. Parisi, R. Zecchina,
  \textit{Analytic and algorithmic solution of random satisfiability problems},
  Science 297 (2002), 812--815.

\bibitem{ART}
  D. Achlioptas, A. Coja-Oghlan,
  \textit{Algorithmic barriers from phase transitions},
  in Proc.\ 49th IEEE Symposium on Foundations of Computer Science (FOCS), 2008, pp.\ 793--802.

\bibitem{DSS15}
  J. Ding, A. Sly, N. Sun,
  \textit{Proof of the satisfiability conjecture for large $k$},
  in Proc.\ 47th Annual ACM Symposium on Theory of Computing (STOC), 2015, pp.\ 59--68.

\bibitem{Ding15}
  J. Ding, A. Sly, N. Sun,
  \textit{Satisfiability threshold for random regular NAE-SAT},
  Communications in Mathematical Physics 341(2) (2016), 435--489.

\bibitem{planted}
  W. Barthel, C. Hartung, M. Leone, F. Ricci-Tersenghi, M. Weigt, R. Zecchina,
  \textit{Hiding solutions in random satisfiability problems: a statistical mechanics approach},
  Physical Review Letters 88 (2002), 188701.

\bibitem{BPR97}
  P. Beame, R. Impagliazzo, J. Kraj\'i\v{c}ek, T. Pitassi, P. Pudl\'ak,
  \textit{Lower bounds on Hilbert's Nullstellensatz and propositional proofs},
  Proc.\ London Mathematical Society 73(3) (1996), 1--26.

\bibitem{GCT}
  K.D. Mulmuley, M. Sohoni,
  \textit{Geometric complexity theory I: An approach to the P vs.\ NP and related problems},
  SIAM Journal on Computing 31(2) (2001), 496--526.

\bibitem{CF-arch}
  B. Sanders,
  \textit{The Coherence Field: TIG Architecture and Operator Calculus},
  Coherence Field Whitepaper 1, 7Site LLC, February 2026.

\end{thebibliography}

\end{document}
